% Options for packages loaded elsewhere
\PassOptionsToPackage{unicode}{hyperref}
\PassOptionsToPackage{hyphens}{url}
\PassOptionsToPackage{dvipsnames,svgnames,x11names}{xcolor}
\documentclass[
  10.5pt,
  a4paper,
]{article}
\usepackage{xcolor}
\usepackage[margin=2.3cm]{geometry}
\usepackage{amsmath,amssymb}
\setcounter{secnumdepth}{-\maxdimen} % remove section numbering
\usepackage{iftex}
\ifPDFTeX
  \usepackage[T1]{fontenc}
  \usepackage[utf8]{inputenc}
  \usepackage{textcomp} % provide euro and other symbols
\else % if luatex or xetex
  \usepackage{unicode-math} % this also loads fontspec
  \defaultfontfeatures{Scale=MatchLowercase}
  \defaultfontfeatures[\rmfamily]{Ligatures=TeX,Scale=1}
\fi
\usepackage{lmodern}
\ifPDFTeX\else
  % xetex/luatex font selection
  \setmainfont[]{Latin Modern Roman}
  \setmonofont[Scale=0.92]{Latin Modern Mono}
\fi
% Use upquote if available, for straight quotes in verbatim environments
\IfFileExists{upquote.sty}{\usepackage{upquote}}{}
\IfFileExists{microtype.sty}{% use microtype if available
  \usepackage[]{microtype}
  \UseMicrotypeSet[protrusion]{basicmath} % disable protrusion for tt fonts
}{}
\makeatletter
\@ifundefined{KOMAClassName}{% if non-KOMA class
  \IfFileExists{parskip.sty}{%
    \usepackage{parskip}
  }{% else
    \setlength{\parindent}{0pt}
    \setlength{\parskip}{6pt plus 2pt minus 1pt}}
}{% if KOMA class
  \KOMAoptions{parskip=half}}
\makeatother
\usepackage{color}
\usepackage{fancyvrb}
\newcommand{\VerbBar}{|}
\newcommand{\VERB}{\Verb[commandchars=\\\{\}]}
\DefineVerbatimEnvironment{Highlighting}{Verbatim}{commandchars=\\\{\}}
% Add ',fontsize=\small' for more characters per line
\newenvironment{Shaded}{}{}
\newcommand{\AlertTok}[1]{\textcolor[rgb]{1.00,0.00,0.00}{\textbf{#1}}}
\newcommand{\AnnotationTok}[1]{\textcolor[rgb]{0.38,0.63,0.69}{\textbf{\textit{#1}}}}
\newcommand{\AttributeTok}[1]{\textcolor[rgb]{0.49,0.56,0.16}{#1}}
\newcommand{\BaseNTok}[1]{\textcolor[rgb]{0.25,0.63,0.44}{#1}}
\newcommand{\BuiltInTok}[1]{\textcolor[rgb]{0.00,0.50,0.00}{#1}}
\newcommand{\CharTok}[1]{\textcolor[rgb]{0.25,0.44,0.63}{#1}}
\newcommand{\CommentTok}[1]{\textcolor[rgb]{0.38,0.63,0.69}{\textit{#1}}}
\newcommand{\CommentVarTok}[1]{\textcolor[rgb]{0.38,0.63,0.69}{\textbf{\textit{#1}}}}
\newcommand{\ConstantTok}[1]{\textcolor[rgb]{0.53,0.00,0.00}{#1}}
\newcommand{\ControlFlowTok}[1]{\textcolor[rgb]{0.00,0.44,0.13}{\textbf{#1}}}
\newcommand{\DataTypeTok}[1]{\textcolor[rgb]{0.56,0.13,0.00}{#1}}
\newcommand{\DecValTok}[1]{\textcolor[rgb]{0.25,0.63,0.44}{#1}}
\newcommand{\DocumentationTok}[1]{\textcolor[rgb]{0.73,0.13,0.13}{\textit{#1}}}
\newcommand{\ErrorTok}[1]{\textcolor[rgb]{1.00,0.00,0.00}{\textbf{#1}}}
\newcommand{\ExtensionTok}[1]{#1}
\newcommand{\FloatTok}[1]{\textcolor[rgb]{0.25,0.63,0.44}{#1}}
\newcommand{\FunctionTok}[1]{\textcolor[rgb]{0.02,0.16,0.49}{#1}}
\newcommand{\ImportTok}[1]{\textcolor[rgb]{0.00,0.50,0.00}{\textbf{#1}}}
\newcommand{\InformationTok}[1]{\textcolor[rgb]{0.38,0.63,0.69}{\textbf{\textit{#1}}}}
\newcommand{\KeywordTok}[1]{\textcolor[rgb]{0.00,0.44,0.13}{\textbf{#1}}}
\newcommand{\NormalTok}[1]{#1}
\newcommand{\OperatorTok}[1]{\textcolor[rgb]{0.40,0.40,0.40}{#1}}
\newcommand{\OtherTok}[1]{\textcolor[rgb]{0.00,0.44,0.13}{#1}}
\newcommand{\PreprocessorTok}[1]{\textcolor[rgb]{0.74,0.48,0.00}{#1}}
\newcommand{\RegionMarkerTok}[1]{#1}
\newcommand{\SpecialCharTok}[1]{\textcolor[rgb]{0.25,0.44,0.63}{#1}}
\newcommand{\SpecialStringTok}[1]{\textcolor[rgb]{0.73,0.40,0.53}{#1}}
\newcommand{\StringTok}[1]{\textcolor[rgb]{0.25,0.44,0.63}{#1}}
\newcommand{\VariableTok}[1]{\textcolor[rgb]{0.10,0.09,0.49}{#1}}
\newcommand{\VerbatimStringTok}[1]{\textcolor[rgb]{0.25,0.44,0.63}{#1}}
\newcommand{\WarningTok}[1]{\textcolor[rgb]{0.38,0.63,0.69}{\textbf{\textit{#1}}}}
\usepackage{longtable,booktabs,array}
\usepackage{caption}
\captionsetup[table]{skip=6pt}
\newcounter{none} % for unnumbered tables
\usepackage{calc} % for calculating minipage widths
% Correct order of tables after \paragraph or \subparagraph
\usepackage{etoolbox}
\makeatletter
\patchcmd\longtable{\par}{\if@noskipsec\mbox{}\fi\par}{}{}
\makeatother
% Allow footnotes in longtable head/foot
\IfFileExists{footnotehyper.sty}{\usepackage{footnotehyper}}{\usepackage{footnote}}
\makesavenoteenv{longtable}
\setlength{\emergencystretch}{3em} % prevent overfull lines
\providecommand{\tightlist}{%
  \setlength{\itemsep}{0pt}\setlength{\parskip}{0pt}}

\usepackage{newunicodechar}
\newunicodechar{}{\ensuremath{\hat{\gamma}}}
\newunicodechar{γ}{\ensuremath{\gamma}}
\newunicodechar{δ}{\ensuremath{\delta}}
\newunicodechar{β}{\ensuremath{\beta}}
\newunicodechar{η}{\ensuremath{\eta}}
\newunicodechar{ε}{\ensuremath{\varepsilon}}
\newunicodechar{₀}{\ensuremath{{}_0}}
\newunicodechar{₁}{\ensuremath{{}_1}}
\newunicodechar{₂}{\ensuremath{{}_2}}
\usepackage{fvextra}
\fvset{breaklines=true, breakanywhere=true, fontsize=\footnotesize}
\DefineVerbatimEnvironment{Highlighting}{Verbatim}{breaklines,breakanywhere,commandchars=\\\{\}}
\RecustomVerbatimEnvironment{verbatim}{Verbatim}{breaklines,breakanywhere}
\usepackage{booktabs}
\usepackage{longtable}
\setlength{\emergencystretch}{3em}
\renewcommand{\arraystretch}{1.15}
\setlength{\tabcolsep}{4pt}
\usepackage{etoolbox}
\AtBeginEnvironment{longtable}{\footnotesize}
\AtBeginEnvironment{quote}{\small}
\usepackage{fancyhdr}
\pagestyle{fancy}
\fancyhf{}
\fancyhead[L]{\small pyxtdpthresh user guide}
\fancyfoot[C]{\small\thepage}
\renewcommand{\headrulewidth}{0.3pt}
\usepackage{bookmark}
\IfFileExists{xurl.sty}{\usepackage{xurl}}{} % add URL line breaks if available
\urlstyle{same}
% fallback for those not using the hyperref driver hyperxmp:
\makeatletter
\@ifundefined{xmpquote}{\newcommand{\xmpquote}[1]{#1}}{}
\makeatother
\hypersetup{
  pdftitle={pyxtdpthresh user guide},
  colorlinks=true,
  linkcolor={black},
  filecolor={Maroon},
  citecolor={Blue},
  urlcolor={blue},
  pdfcreator={LaTeX via pandoc}}

\title{pyxtdpthresh user guide}
\usepackage{etoolbox}
\makeatletter
\providecommand{\subtitle}[1]{% add subtitle to \maketitle
  \apptocmd{\@title}{\par {\large #1 \par}}{}{}
}
\makeatother
\subtitle{Version 0.0.1 (alpha)}
\author{}
\date{03 November 2026}

\begin{document}
\maketitle

{
\setcounter{tocdepth}{2}
\tableofcontents
}
This guide explains every argument of the two estimators, how each one
corresponds to the Stata option it ports, and how to read what comes
back. It assumes you know what a dynamic panel threshold model is; the
equations and the algorithms are in the methodological note, Appendix A
of the PDF edition of this guide (\texttt{user\_guide.pdf}, built from
\texttt{docs/methodological\_note.tex}). The Stata help files
(\texttt{xtdpthresh.sthlp}, \texttt{xthenreg.sthlp}) remain the
reference for the commands; this guide says where the Python package
differs from them.

The code blocks of this guide run top to bottom as one script
(\texttt{tests/integration/test\_examples.py} runs them), so a block may
use the names defined in the blocks before it. A block whose first line
is \texttt{\#\ not\ run\ by\ the\ tests} is shown for illustration only.

\begin{center}\rule{0.5\linewidth}{0.5pt}\end{center}

\section{1. Before you start}\label{before-you-start}

\textbf{Data.} Both estimators take a pandas \texttt{DataFrame} in long
form, one row per (unit, period). Name the unit and time columns with
\texttt{id=} and \texttt{time=}, or give the frame a two-level
\texttt{MultiIndex} (unit, time) and leave both out.
\texttt{pyx.from\_stata("file.dta")} reads a \texttt{.dta} file and uses
the \texttt{xtset} variables Stata recorded in it; a file saved without
\texttt{xtset} is returned as it is and you pass \texttt{id=} and
\texttt{time=} yourself (the names are never guessed).

\texttt{xtdpthresh} accepts unbalanced panels. \texttt{xthenreg} needs a
strongly balanced panel (every unit observed in every period) and stops
with an error otherwise, as the Stata command does.

\textbf{Time-series operators.} Variable names in every role may carry
Stata\textquotesingle s operators as text: \texttt{"L.x"},
\texttt{"L2.x"}, \texttt{"D.x"}, \texttt{"S12.x"}, \texttt{"L(1/2).x"}.
They are computed on the calendar, as in Stata: \texttt{L.x} is the
value of the \emph{period} before, so a gap in a unit\textquotesingle s
dates gives a missing value. The dependent variable may not carry an
operator. The rules of the help file apply: each variable sits in only
one role, a lag of the dependent variable is declared
\texttt{predetermined}, and differences and leads of the dependent
variable are refused.

\textbf{The example used below.} Hansen\textquotesingle s (1999) firm
investment panel, 565 firms observed from 1973 to 1987, with the five
extreme firms removed as in Seo and Shin (2016).

\begin{Shaded}
\begin{Highlighting}[]
\ImportTok{import}\NormalTok{ pyxtdpthresh }\ImportTok{as}\NormalTok{ pyx}

\NormalTok{df }\OperatorTok{=}\NormalTok{ pyx.load\_dataset(}\StringTok{"invest"}\NormalTok{)}
\NormalTok{df }\OperatorTok{=}\NormalTok{ df[}\OperatorTok{\textasciitilde{}}\NormalTok{df[}\StringTok{"firm"}\NormalTok{].isin([}\DecValTok{137}\NormalTok{, }\DecValTok{391}\NormalTok{, }\DecValTok{407}\NormalTok{, }\DecValTok{488}\NormalTok{, }\DecValTok{538}\NormalTok{])]}

\CommentTok{\# The model of Seo and Shin (2016, section 7): all four slopes change at a leverage threshold,}
\CommentTok{\# cash flow, Tobin\textquotesingle{}s q and leverage are endogenous, instruments are lags 3 to 8, collapsed.}
\NormalTok{MODEL }\OperatorTok{=} \BuiltInTok{dict}\NormalTok{(}
\NormalTok{    q}\OperatorTok{=}\StringTok{"debt"}\NormalTok{, endogenous}\OperatorTok{=}\NormalTok{[}\StringTok{"cashflow"}\NormalTok{, }\StringTok{"tobin\_q"}\NormalTok{, }\StringTok{"debt"}\NormalTok{],}
\NormalTok{    maxlag}\OperatorTok{=}\NormalTok{(}\DecValTok{3}\NormalTok{, }\DecValTok{8}\NormalTok{), collapse}\OperatorTok{=}\VariableTok{True}\NormalTok{, }\BuiltInTok{id}\OperatorTok{=}\StringTok{"firm"}\NormalTok{, time}\OperatorTok{=}\StringTok{"year"}\NormalTok{,}
\NormalTok{)}
\end{Highlighting}
\end{Shaded}

\textbf{Other ways in.} The unit and time columns can be the two levels
of the index; \texttt{pyx.list\_datasets()} names the bundled data
(\texttt{"invest"} today); \texttt{pyx.from\_stata} reads a
\texttt{.dta} file.

\begin{Shaded}
\begin{Highlighting}[]
\NormalTok{indexed }\OperatorTok{=}\NormalTok{ df.set\_index([}\StringTok{"firm"}\NormalTok{, }\StringTok{"year"}\NormalTok{])}
\NormalTok{no\_names }\OperatorTok{=}\NormalTok{ \{k: v }\ControlFlowTok{for}\NormalTok{ k, v }\KeywordTok{in}\NormalTok{ MODEL.items() }\ControlFlowTok{if}\NormalTok{ k }\KeywordTok{not} \KeywordTok{in}\NormalTok{ (}\StringTok{"id"}\NormalTok{, }\StringTok{"time"}\NormalTok{)\}}
\NormalTok{same }\OperatorTok{=}\NormalTok{ pyx.xtdpthresh(indexed, }\StringTok{"invest"}\NormalTok{, bootstrap}\OperatorTok{=}\VariableTok{False}\NormalTok{, }\OperatorTok{**}\NormalTok{no\_names)}
\BuiltInTok{print}\NormalTok{(}\BuiltInTok{round}\NormalTok{(same.gamma, }\DecValTok{4}\NormalTok{))}

\BuiltInTok{print}\NormalTok{(pyx.list\_datasets())        }\CommentTok{\# name {-}\textgreater{} description}

\CommentTok{\# A .dta saved without xtset gives a plain frame: the names are never guessed.}
\NormalTok{df[[}\StringTok{"firm"}\NormalTok{, }\StringTok{"year"}\NormalTok{, }\StringTok{"invest"}\NormalTok{, }\StringTok{"debt"}\NormalTok{]].to\_stata(}\StringTok{"panel.dta"}\NormalTok{, write\_index}\OperatorTok{=}\VariableTok{False}\NormalTok{)}
\NormalTok{raw }\OperatorTok{=}\NormalTok{ pyx.from\_stata(}\StringTok{"panel.dta"}\NormalTok{)}
\BuiltInTok{print}\NormalTok{(raw.attrs[}\StringTok{"id"}\NormalTok{], raw.attrs[}\StringTok{"time"}\NormalTok{])      }\CommentTok{\# None None: pass id= and time= yourself}
\end{Highlighting}
\end{Shaded}

A file saved after \texttt{xtset} (or \texttt{xtset\ …,\ delta()}) gives
a frame whose index levels are the panel and time variables Stata
recorded, with their names in \texttt{attrs{[}"id"{]}},
\texttt{attrs{[}"time"{]}} and the time unit in
\texttt{attrs{[}"delta"{]}}; \texttt{id=} and \texttt{time=} can then be
left out. Stata value labels are not applied.

\begin{center}\rule{0.5\linewidth}{0.5pt}\end{center}

\section{\texorpdfstring{2. Calling
\texttt{xtdpthresh}}{2. Calling xtdpthresh}}\label{calling-xtdpthresh}

\begin{Shaded}
\begin{Highlighting}[]
\NormalTok{res = pyx.xtdpthresh(data, y, x=(), *, q, id=None, time=None, where=None, grid\_values=None, **options)}
\end{Highlighting}
\end{Shaded}

{\def\LTcaptype{none} % do not increment counter
\begin{longtable}[]{@{}
  >{\raggedright\arraybackslash}p{(\linewidth - 4\tabcolsep) * \real{0.1414}}
  >{\raggedright\arraybackslash}p{(\linewidth - 4\tabcolsep) * \real{0.2323}}
  >{\raggedright\arraybackslash}p{(\linewidth - 4\tabcolsep) * \real{0.6263}}@{}}
\toprule\noalign{}
\begin{minipage}[b]{\linewidth}\raggedright
Argument
\end{minipage} & \begin{minipage}[b]{\linewidth}\raggedright
Stata
\end{minipage} & \begin{minipage}[b]{\linewidth}\raggedright
Meaning
\end{minipage} \\
\midrule\noalign{}
\endhead
\bottomrule\noalign{}
\endlastfoot
\texttt{data} & the data in memory & long \texttt{DataFrame} \\
\texttt{y} & \texttt{depvar} & dependent variable (a plain column
name) \\
\texttt{x} & \texttt{indepvars} & strictly exogenous regressors, a name
or a list \\
\texttt{q} & \texttt{qx(varname)} & threshold variable,
\textbf{required} \\
\texttt{endogenous} & \texttt{endogenous(varlist)} & regressors
correlated with the current error \\
\texttt{predetermined} & \texttt{predetermined(varlist)} & regressors
uncorrelated with current and later errors \\
\texttt{exogenous} & \texttt{exogenous(varlist)} & more strictly
exogenous regressors, treated like \texttt{x} \\
\texttt{iv} & \texttt{iv(varlist)} & external instruments, entered with
their period-t values \\
\texttt{id}, \texttt{time} & \texttt{xtset\ id\ time} & unit and time
columns (or index level names) \\
\texttt{where} & \texttt{if} & condition selecting the equations (see
below) \\
\texttt{grid\_values} & (none) & candidate thresholds to use instead of
the grid built from \texttt{grid} \\
\texttt{**options} & the other options & Section 3 \\
\end{longtable}
}

\texttt{endogenous}, \texttt{predetermined}, \texttt{exogenous} and
\texttt{iv} are keyword arguments given as a name or a list of names.
Anything else the call receives is an option of Section 3; an unknown
name raises a \texttt{TypeError}
(\texttt{…\ got\ an\ unexpected\ keyword\ argument\ \textquotesingle{}gridd\textquotesingle{}}).

The model is
\texttt{y\ =\ x\textquotesingle{}β\ +\ (1,\ x\textquotesingle{})δ\ ·\ 1(q\ \textgreater{}\ γ)\ +\ η\ +\ ε}.
In the dynamic model (the default) the lag of \texttt{y} is added to
\texttt{x} automatically; \texttt{static=True} removes it. The slopes
are \texttt{β} below the threshold and \texttt{β\ +\ δ₂} above it;
\texttt{δ₁} is the shift of the intercept.

\textbf{How the roles are instrumented.} With \texttt{method="fd"}, an
endogenous regressor is instrumented by its levels from t-2, a
predetermined one from t-1, a strictly exogenous one by its own
transformed value; with \texttt{method="fod"} the dates are t-1 and t.
\texttt{maxlag} and \texttt{collapse} shape these sets, \texttt{iv} adds
external ones.

\textbf{\texttt{where} and \texttt{in}.} \texttt{where} is a string
evaluated with \texttt{DataFrame.eval}
(\texttt{"year\ \textgreater{}=\ 1980"}, \texttt{"firm\ !=\ 12"}) or a
boolean array of the length of the frame. Two differences from
Stata\textquotesingle s \texttt{if}: a missing value in the condition
counts as \emph{false} (Stata counts it as true; the rows with a missing
variable are removed afterwards anyway), and there is no \texttt{in};
select rows of the frame instead. By default (\texttt{history="panel"})
the condition restricts the \emph{equations}, and observations outside
it remain available as instruments.

\begin{Shaded}
\begin{Highlighting}[]
\CommentTok{\# Point estimates only: the bootstrap is the slow part, bootstrap=False skips it.}
\NormalTok{fast }\OperatorTok{=}\NormalTok{ pyx.xtdpthresh(df, }\StringTok{"invest"}\NormalTok{, bootstrap}\OperatorTok{=}\VariableTok{False}\NormalTok{, }\OperatorTok{**}\NormalTok{MODEL)}
\BuiltInTok{print}\NormalTok{(}\BuiltInTok{round}\NormalTok{(fast.gamma, }\DecValTok{4}\NormalTok{), fast.n\_obs, fast.n\_units, fast.n\_instruments)}

\CommentTok{\# Equations from 1980 on; the earlier observations still serve as instruments.}
\NormalTok{late }\OperatorTok{=}\NormalTok{ pyx.xtdpthresh(df, }\StringTok{"invest"}\NormalTok{, where}\OperatorTok{=}\StringTok{"year \textgreater{}= 1980"}\NormalTok{, bootstrap}\OperatorTok{=}\VariableTok{False}\NormalTok{, }\OperatorTok{**}\NormalTok{MODEL)}
\BuiltInTok{print}\NormalTok{(late.n\_obs)}

\CommentTok{\# Your own candidate thresholds, instead of the grid built from \textasciigrave{}grid\textasciigrave{}.}
\ImportTok{import}\NormalTok{ numpy }\ImportTok{as}\NormalTok{ np}

\NormalTok{own }\OperatorTok{=}\NormalTok{ pyx.xtdpthresh(df, }\StringTok{"invest"}\NormalTok{, grid\_values}\OperatorTok{=}\NormalTok{np.linspace(}\FloatTok{0.05}\NormalTok{, }\FloatTok{0.55}\NormalTok{, }\DecValTok{11}\NormalTok{), bootstrap}\OperatorTok{=}\VariableTok{False}\NormalTok{, }\OperatorTok{**}\NormalTok{MODEL)}
\BuiltInTok{print}\NormalTok{(}\BuiltInTok{round}\NormalTok{(own.gamma, }\DecValTok{4}\NormalTok{))}
\end{Highlighting}
\end{Shaded}

A complete call, with inference, a seed, and the test of the threshold
estimated by Hansen (1999), 0.0157 (this is the second example of the
Stata help file):

\begin{Shaded}
\begin{Highlighting}[]
\NormalTok{res }\OperatorTok{=}\NormalTok{ pyx.xtdpthresh(df, }\StringTok{"invest"}\NormalTok{, gamma\_test}\OperatorTok{=}\FloatTok{0.0157}\NormalTok{, seed}\OperatorTok{=}\DecValTok{12345}\NormalTok{, }\OperatorTok{**}\NormalTok{MODEL)}
\BuiltInTok{print}\NormalTok{(res)}
\end{Highlighting}
\end{Shaded}

\textbf{External instruments.} \texttt{iv} takes names (with operators)
of variables that are not regressors:

\begin{Shaded}
\begin{Highlighting}[]
\NormalTok{df[}\StringTok{"peer\_q"}\NormalTok{] }\OperatorTok{=}\NormalTok{ df.groupby(}\StringTok{"year"}\NormalTok{)[}\StringTok{"tobin\_q"}\NormalTok{].transform(}\StringTok{"mean"}\NormalTok{)   }\CommentTok{\# yearly mean of Tobin\textquotesingle{}s q}
\NormalTok{withiv }\OperatorTok{=}\NormalTok{ pyx.xtdpthresh(df, }\StringTok{"invest"}\NormalTok{, iv}\OperatorTok{=}\NormalTok{[}\StringTok{"L2.peer\_q"}\NormalTok{], bootstrap}\OperatorTok{=}\VariableTok{False}\NormalTok{, }\OperatorTok{**}\NormalTok{MODEL)}
\BuiltInTok{print}\NormalTok{(withiv.n\_instruments)}
\end{Highlighting}
\end{Shaded}

A variable that takes the same value for every unit in a period (such as
\texttt{peer\_q}) is identified through the period constants: without
\texttt{collapse}, its own columns repeat the constants and are dropped,
and \texttt{res.scalars{[}"N\_iv\_common"{]}} counts them.

\begin{center}\rule{0.5\linewidth}{0.5pt}\end{center}

\section{\texorpdfstring{3. Options of
\texttt{xtdpthresh}}{3. Options of xtdpthresh}}\label{options-of-xtdpthresh}

Python names are in \texttt{snake\_case}; the Stata name is given for
each. Defaults are those of the Stata command. Every rule of the Stata
command is checked in the same order and raises
\texttt{SpecificationError} (Section 10) carrying
Stata\textquotesingle s return code.

\subsection{Model}\label{model}

{\def\LTcaptype{none} % do not increment counter
\begin{longtable}[]{@{}
  >{\raggedright\arraybackslash}p{(\linewidth - 6\tabcolsep) * \real{0.1188}}
  >{\raggedright\arraybackslash}p{(\linewidth - 6\tabcolsep) * \real{0.2079}}
  >{\raggedright\arraybackslash}p{(\linewidth - 6\tabcolsep) * \real{0.0792}}
  >{\raggedright\arraybackslash}p{(\linewidth - 6\tabcolsep) * \real{0.5941}}@{}}
\toprule\noalign{}
\begin{minipage}[b]{\linewidth}\raggedright
Python
\end{minipage} & \begin{minipage}[b]{\linewidth}\raggedright
Stata
\end{minipage} & \begin{minipage}[b]{\linewidth}\raggedright
Default
\end{minipage} & \begin{minipage}[b]{\linewidth}\raggedright
Notes
\end{minipage} \\
\midrule\noalign{}
\endhead
\bottomrule\noalign{}
\endlastfoot
\texttt{kink} & \texttt{kink} & \texttt{False} & Continuous effect at
the threshold instead of a jump. Put \texttt{q} in \texttt{x},
\texttt{endogenous} or \texttt{predetermined} so that it has a slope
below γ; otherwise a warning says the model is a one-sided hinge. \\
\texttt{static} & \texttt{static} & \texttt{False} & No automatic lag of
\texttt{y}. Accepted by the Stata command; the help file of version
0.9.38 does not describe it. \\
\texttt{time\_effects} & \texttt{td} & \texttt{False} & Partial common
time effects out of the equations (dummies are not entered by hand). \\
\texttt{history} &
\texttt{history(panel\textbackslash{}\textbar{}sample)} &
\texttt{"panel"} & With \texttt{"sample"}, \texttt{where} also bounds
the lag history; then no operated term (\texttt{L.x}) may be used. Build
the lags as columns first. \\
\end{longtable}
}

\subsection{Instruments}\label{instruments}

{\def\LTcaptype{none} % do not increment counter
\begin{longtable}[]{@{}
  >{\raggedright\arraybackslash}p{(\linewidth - 6\tabcolsep) * \real{0.1200}}
  >{\raggedright\arraybackslash}p{(\linewidth - 6\tabcolsep) * \real{0.1600}}
  >{\raggedright\arraybackslash}p{(\linewidth - 6\tabcolsep) * \real{0.0700}}
  >{\raggedright\arraybackslash}p{(\linewidth - 6\tabcolsep) * \real{0.6500}}@{}}
\toprule\noalign{}
\begin{minipage}[b]{\linewidth}\raggedright
Python
\end{minipage} & \begin{minipage}[b]{\linewidth}\raggedright
Stata
\end{minipage} & \begin{minipage}[b]{\linewidth}\raggedright
Default
\end{minipage} & \begin{minipage}[b]{\linewidth}\raggedright
Notes
\end{minipage} \\
\midrule\noalign{}
\endhead
\bottomrule\noalign{}
\endlastfoot
\texttt{maxlag} & \texttt{maxlag(\#\ {[}\#{]})} & all lags & An integer
\texttt{L} uses lags 1 to \texttt{L} for predetermined variables and 2
to \texttt{L} for endogenous variables and \texttt{y}; a pair
\texttt{(a,\ b)} uses \texttt{a} to \texttt{b} with the same minimums.
\texttt{None} means no limit. Without it the number of instruments grows
with the square of T. \\
\texttt{collapse} & \texttt{collapse} & \texttt{False} & One instrument
column per lag instead of one per lag and period (Roodman 2009). Fewer
instruments, possibly weaker identification of γ. \\
\texttt{iv\_collapse} & \texttt{iv(…,\ collapse)} & \texttt{False} &
Collapse only the external instruments. \texttt{collapse=True} implies
it. \\
\end{longtable}
}

\subsection{Transformation and threshold
search}\label{transformation-and-threshold-search}

{\def\LTcaptype{none} % do not increment counter
\begin{longtable}[]{@{}
  >{\raggedright\arraybackslash}p{(\linewidth - 6\tabcolsep) * \real{0.1000}}
  >{\raggedright\arraybackslash}p{(\linewidth - 6\tabcolsep) * \real{0.2600}}
  >{\raggedright\arraybackslash}p{(\linewidth - 6\tabcolsep) * \real{0.1000}}
  >{\raggedright\arraybackslash}p{(\linewidth - 6\tabcolsep) * \real{0.5400}}@{}}
\toprule\noalign{}
\begin{minipage}[b]{\linewidth}\raggedright
Python
\end{minipage} & \begin{minipage}[b]{\linewidth}\raggedright
Stata
\end{minipage} & \begin{minipage}[b]{\linewidth}\raggedright
Default
\end{minipage} & \begin{minipage}[b]{\linewidth}\raggedright
Notes
\end{minipage} \\
\midrule\noalign{}
\endhead
\bottomrule\noalign{}
\endlastfoot
\texttt{method} & \texttt{method(fd\textbackslash{}\textbar{}fod)} &
\texttt{"fd"} & First differences, or forward orthogonal deviations. FOD
discards fewer equations when the panel has interior gaps (an interior
gap removes three FD equations and two FOD equations). \\
\texttt{grid} & \texttt{grid(\#)} & \texttt{100} & Points of the
threshold search, at least 10. Uniform grids of 100, 199 and 397 points
are nested, so a finer grid checks the sensitivity of . \\
\texttt{grid\_type} &
\texttt{gridtype(uniform\textbackslash{}\textbar{}quantile)} &
\texttt{"uniform"} & Equally spaced in \texttt{q}, or at its empirical
quantiles; used for the estimation grid and the confidence-set grid. \\
\texttt{trim} & \texttt{trim(\#)} & \texttt{0.10} & Total trimming rate
of \texttt{q}, split between the tails; 0.10 restricts the grid to the
5th--95th percentiles. Between 0.01 and 0.45. (\texttt{xthenreg} uses
0.4.) \\
\texttt{grid\_sample} &
\texttt{gridsample(effective\textbackslash{}\textbar{}observed)} &
\texttt{"effective"} & Which values of \texttt{q} define the trimming
bounds and the regime counts of \texttt{min\_regime}: those that enter
the transformed equations, or those in the retained rows. \\
\texttt{min\_regime} & \texttt{minregime(\#)} & \texttt{0} & Raise the
minimum number of observations in each regime. The default, and the
least allowed, is the size of the smaller trimmed tail. \\
\texttt{refine} & \texttt{refine(\#)} & \texttt{0} & Rounds (0--20) of
local refinement after the grid search: the observed values of
\texttt{q} next to  become candidates. The confidence set and tests use
the initial grid only. Not allowed with \texttt{kink}, whose  is
refined continuously. \\
\texttt{grid\_values} & (none) & \texttt{None} & An argument of the
call, not an option: an explicit array of candidates. \\
\end{longtable}
}

\subsection{Inference}\label{inference}

{\def\LTcaptype{none} % do not increment counter
\begin{longtable}[]{@{}
  >{\raggedright\arraybackslash}p{(\linewidth - 6\tabcolsep) * \real{0.1287}}
  >{\raggedright\arraybackslash}p{(\linewidth - 6\tabcolsep) * \real{0.2574}}
  >{\raggedright\arraybackslash}p{(\linewidth - 6\tabcolsep) * \real{0.0990}}
  >{\raggedright\arraybackslash}p{(\linewidth - 6\tabcolsep) * \real{0.5149}}@{}}
\toprule\noalign{}
\begin{minipage}[b]{\linewidth}\raggedright
Python
\end{minipage} & \begin{minipage}[b]{\linewidth}\raggedright
Stata
\end{minipage} & \begin{minipage}[b]{\linewidth}\raggedright
Default
\end{minipage} & \begin{minipage}[b]{\linewidth}\raggedright
Notes
\end{minipage} \\
\midrule\noalign{}
\endhead
\bottomrule\noalign{}
\endlastfoot
\texttt{bootstrap} & \texttt{noboot} (negated) & \texttt{True} &
\texttt{False} skips the confidence set and the linearity test. Point
estimates, Hansen and AR tests remain. \\
\texttt{n\_boot} & \texttt{boot(\#)} & \texttt{299} & Bootstrap
replications, at least 10. With B draws the smallest possible p-value is
1/(B + 1). \\
\texttt{grid\_ci} & \texttt{gridci(\#)} & \texttt{100} & Candidate
thresholds of the confidence-set grid, at least 10. \\
\texttt{boot\_type} &
\texttt{boottype(wild\textbackslash{}\textbar{}unit)} & \texttt{"wild"}
& Scheme of the confidence set and of \texttt{gamma\_test}.
\texttt{"unit"} resamples whole units, follows Algorithm 1 of Gong and
Seo (2026), needs \texttt{method="fd"} without \texttt{kink} or
\texttt{time\_effects}, and is much slower. The linearity test always
uses the wild scheme. \\
\texttt{seed} & \texttt{rseed(\#)} & drawn & Integer in 0 \ldots{}
2³¹-1. When the bootstrap is on and \texttt{seed} is omitted, a seed is
drawn and recorded: \texttt{res.seed} replays the run. Separate seeds
are derived for the confidence set, the linearity test and
\texttt{gamma\_test}, so changing one leaves the draws of the others
unchanged. \\
\texttt{linearity\_test} & \texttt{notest} (negated) & \texttt{True} &
\texttt{False} skips the bootstrap linearity test. \\
\texttt{gamma\_test} & \texttt{citest(\#)} & \texttt{None} & Also run
the bootstrap test of H₀: γ = \#, with a value of \texttt{q}. Needs the
bootstrap. It is evaluated after the other components with its own seed,
so asking for it changes nothing else. The result is in
\texttt{res.scalars{[}"citest\_p"{]}} and
\texttt{{[}"citest\_accept"{]}}. \\
\texttt{continuity\_test} & \texttt{conttest} & \texttt{False} &
Bootstrap test of continuity (is the kink model enough?), result in
\texttt{res.scalars{[}"pval\_cont"{]}}. Needs the bootstrap and the
linearity test, no \texttt{kink}, and \texttt{q} as a contemporaneous
regressor. Accepted by the command; not described in the 0.9.38 help
file. \\
\texttt{vce} & \texttt{vce(robust\textbackslash{}\textbar{}windmeijer)}
& \texttt{"robust"} & Cluster-robust two-step sandwich, or a
Windmeijer-type correction. The variance includes the estimation error
of ; the one conditional on  is \texttt{res.fit.variance.cond}
(\texttt{e(V\_cond)}). \\
\texttt{bw\_scale} & \texttt{bwscale(\#)} & \texttt{1.5} & Scale of the
kernel bandwidth \texttt{h\ =\ \#\ ·\ 1.06\ ·\ s\_q\ ·\ n\^{}(-1/5)} of
the jump model\textquotesingle s joint variance. Only the reported
variance and the AR statistics that include  depend on it. \\
\texttt{center} & \texttt{nocenter} (negated) & \texttt{True} &
\texttt{False} uses the uncentered moment covariance, as
\texttt{xtabond2} does. \\
\texttt{level} & \texttt{level(\#)} & \texttt{95} & Confidence level,
between 10 and 99.99. \\
\texttt{coef\_boot} &
\texttt{coefboot(twostep\textbackslash{}\textbar{}onestep\textbackslash{}\textbar{}none)}
& \texttt{"none"} & Legacy coefficient bootstrap that searches the
threshold again in each draw. A note says so; the intervals are in
\texttt{res.fit.bootstrap.coefficient.interval} (\texttt{e(b\_bootci)})
and the printed tables keep the analytic ones. \texttt{"twostep"} is
refused with \texttt{time\_effects}. \\
\texttt{coef\_ci\_type} &
\texttt{coefcitype(symmetric\textbackslash{}\textbar{}percentile)} &
\texttt{"symmetric"} & Kind of interval of that bootstrap. \\
\end{longtable}
}

\subsection{Reporting}\label{reporting}

{\def\LTcaptype{none} % do not increment counter
\begin{longtable}[]{@{}
  >{\raggedright\arraybackslash}p{(\linewidth - 6\tabcolsep) * \real{0.1100}}
  >{\raggedright\arraybackslash}p{(\linewidth - 6\tabcolsep) * \real{0.1700}}
  >{\raggedright\arraybackslash}p{(\linewidth - 6\tabcolsep) * \real{0.0700}}
  >{\raggedright\arraybackslash}p{(\linewidth - 6\tabcolsep) * \real{0.6500}}@{}}
\toprule\noalign{}
\begin{minipage}[b]{\linewidth}\raggedright
Python
\end{minipage} & \begin{minipage}[b]{\linewidth}\raggedright
Stata
\end{minipage} & \begin{minipage}[b]{\linewidth}\raggedright
Default
\end{minipage} & \begin{minipage}[b]{\linewidth}\raggedright
Notes
\end{minipage} \\
\midrule\noalign{}
\endhead
\bottomrule\noalign{}
\endlastfoot
\texttt{warn} & \texttt{nowarn} (negated) & \texttt{True} &
\texttt{False} silences nonfatal warnings and most notes. Notes on
dropped instrument columns and a missing Hansen J, and the warnings for
an incomplete or empty confidence set, are always shown. \\
\texttt{verbose} & \texttt{verbose} & \texttt{False} & Accepted,
\textbf{no effect} in Python. The per-candidate table of the confidence
set is \texttt{res.ci\_table()}. \\
\texttt{export\_gmm} & \texttt{exportgmm} & \texttt{False} & Accepted,
\textbf{no effect} in Python. \\
\end{longtable}
}

\subsection{Combinations the commands
refuse}\label{combinations-the-commands-refuse}

\begin{itemize}
\tightlist
\item
  \texttt{n\_boot\ \textless{}\ 10} or
  \texttt{grid\_ci\ \textless{}\ 10} with the bootstrap on;
  \texttt{grid\ \textless{}\ 10}.
\item
  \texttt{gamma\_test} without the bootstrap; \texttt{continuity\_test}
  without the bootstrap and the linearity test.
\item
  \texttt{refine\ \textgreater{}\ 0} with \texttt{kink}.
\item
  \texttt{boot\_type="unit"} with \texttt{method="fod"} or with
  \texttt{kink} (checked when the bootstrap is on), and with
  \texttt{time\_effects}.
\item
  \texttt{coef\_boot="twostep"} with \texttt{time\_effects}.
\item
  A \texttt{maxlag} pair with the low end above the high end.
\item
  \texttt{history="sample"} with a lagged or differenced term.
\end{itemize}

\subsection{Retired options}\label{retired-options}

These existed in earlier versions of the Stata command and are refused
with an explanation: \texttt{method="system"},
\texttt{boot\_type="exact"}, \texttt{coef\_boot="gs"}, and the arguments
\texttt{search\_mode}, \texttt{search\_max}, \texttt{search\_tol},
\texttt{lev\_maxlag} and \texttt{iv\_maxlag}.

\subsection{The options in use}\label{the-options-in-use}

Each option below is used once, on the model of Section 1. Point
estimates first (\texttt{bootstrap=False} is instant), then the options
of the bootstrap with a light one (99 draws, 40 candidates).

\begin{Shaded}
\begin{Highlighting}[]
\KeywordTok{def}\NormalTok{ fit(}\OperatorTok{**}\NormalTok{changes):}
    \CommentTok{"""Point estimates of the model of Section 1, with some options changed."""}
    \ControlFlowTok{return}\NormalTok{ pyx.xtdpthresh(df, }\StringTok{"invest"}\NormalTok{, }\OperatorTok{**}\NormalTok{\{}\OperatorTok{**}\NormalTok{MODEL, }\StringTok{"bootstrap"}\NormalTok{: }\VariableTok{False}\NormalTok{, }\OperatorTok{**}\NormalTok{changes\})}


\KeywordTok{def}\NormalTok{ show(label, r):}
    \BuiltInTok{print}\NormalTok{(}\SpecialStringTok{f"}\SpecialCharTok{\{}\NormalTok{label}\SpecialCharTok{:26s\}}\SpecialStringTok{ gamma=}\SpecialCharTok{\{}\NormalTok{r}\SpecialCharTok{.}\NormalTok{gamma}\SpecialCharTok{:.4f\}}\SpecialStringTok{  equations=}\SpecialCharTok{\{}\NormalTok{r}\SpecialCharTok{.}\NormalTok{n\_obs}\SpecialCharTok{\}}\SpecialStringTok{  instruments=}\SpecialCharTok{\{}\NormalTok{r}\SpecialCharTok{.}\NormalTok{n\_instruments}\SpecialCharTok{\}}\SpecialStringTok{"}
          \SpecialStringTok{f"  Hansen p=}\SpecialCharTok{\{}\NormalTok{r}\SpecialCharTok{.}\NormalTok{scalars[}\StringTok{\textquotesingle{}hansen\_p\textquotesingle{}}\NormalTok{]}\SpecialCharTok{:.3f\}}\SpecialStringTok{"}\NormalTok{)}


\CommentTok{\# Model}
\NormalTok{show(}\StringTok{"static (lags 2{-}8)"}\NormalTok{, fit(static}\OperatorTok{=}\VariableTok{True}\NormalTok{, maxlag}\OperatorTok{=}\NormalTok{(}\DecValTok{2}\NormalTok{, }\DecValTok{8}\NormalTok{)))   }\CommentTok{\# no automatic L.invest}
\NormalTok{show(}\StringTok{"time effects (td)"}\NormalTok{, fit(time\_effects}\OperatorTok{=}\VariableTok{True}\NormalTok{))}
\NormalTok{show(}\StringTok{"kink"}\NormalTok{, fit(kink}\OperatorTok{=}\VariableTok{True}\NormalTok{))}
\NormalTok{show(}\StringTok{"history=\textquotesingle{}sample\textquotesingle{}, 1976+"}\NormalTok{, fit(history}\OperatorTok{=}\StringTok{"sample"}\NormalTok{, where}\OperatorTok{=}\StringTok{"year \textgreater{}= 1976"}\NormalTok{))}

\CommentTok{\# Instruments}
\NormalTok{show(}\StringTok{"maxlag (3, 5)"}\NormalTok{, fit(maxlag}\OperatorTok{=}\NormalTok{(}\DecValTok{3}\NormalTok{, }\DecValTok{5}\NormalTok{)))}
\NormalTok{show(}\StringTok{"no collapse, lags 3{-}5"}\NormalTok{, fit(collapse}\OperatorTok{=}\VariableTok{False}\NormalTok{, maxlag}\OperatorTok{=}\NormalTok{(}\DecValTok{3}\NormalTok{, }\DecValTok{5}\NormalTok{)))}
\NormalTok{show(}\StringTok{"external, iv\_collapse"}\NormalTok{, fit(iv}\OperatorTok{=}\NormalTok{[}\StringTok{"L2.cashflow"}\NormalTok{], iv\_collapse}\OperatorTok{=}\VariableTok{True}\NormalTok{, collapse}\OperatorTok{=}\VariableTok{False}\NormalTok{, maxlag}\OperatorTok{=}\NormalTok{(}\DecValTok{3}\NormalTok{, }\DecValTok{5}\NormalTok{)))}

\CommentTok{\# Transformation, search, variance}
\NormalTok{show(}\StringTok{"fod"}\NormalTok{, fit(method}\OperatorTok{=}\StringTok{"fod"}\NormalTok{))}
\NormalTok{show(}\StringTok{"trim 0.20"}\NormalTok{, fit(trim}\OperatorTok{=}\FloatTok{0.20}\NormalTok{))}
\NormalTok{show(}\StringTok{"quantile grid"}\NormalTok{, fit(grid\_type}\OperatorTok{=}\StringTok{"quantile"}\NormalTok{))}
\NormalTok{show(}\StringTok{"grid\_sample=\textquotesingle{}observed\textquotesingle{}"}\NormalTok{, fit(grid\_sample}\OperatorTok{=}\StringTok{"observed"}\NormalTok{))}
\NormalTok{show(}\StringTok{"min\_regime 500"}\NormalTok{, fit(min\_regime}\OperatorTok{=}\DecValTok{500}\NormalTok{))}
\NormalTok{show(}\StringTok{"refine 3"}\NormalTok{, fit(refine}\OperatorTok{=}\DecValTok{3}\NormalTok{))}
\NormalTok{show(}\StringTok{"windmeijer"}\NormalTok{, fit(vce}\OperatorTok{=}\StringTok{"windmeijer"}\NormalTok{))}
\NormalTok{show(}\StringTok{"uncentered covariance"}\NormalTok{, fit(center}\OperatorTok{=}\VariableTok{False}\NormalTok{))}
\NormalTok{show(}\StringTok{"bw\_scale 3"}\NormalTok{, fit(bw\_scale}\OperatorTok{=}\DecValTok{3}\NormalTok{))}
\NormalTok{show(}\StringTok{"warn=False"}\NormalTok{, fit(warn}\OperatorTok{=}\VariableTok{False}\NormalTok{))}
\end{Highlighting}
\end{Shaded}

\begin{Shaded}
\begin{Highlighting}[]
\NormalTok{infer }\OperatorTok{=} \BuiltInTok{dict}\NormalTok{(n\_boot}\OperatorTok{=}\DecValTok{99}\NormalTok{, grid\_ci}\OperatorTok{=}\DecValTok{40}\NormalTok{, seed}\OperatorTok{=}\DecValTok{1}\NormalTok{)     }\CommentTok{\# a light bootstrap, for illustration}


\KeywordTok{def}\NormalTok{ full(}\OperatorTok{**}\NormalTok{changes):}
    \ControlFlowTok{return}\NormalTok{ pyx.xtdpthresh(df, }\StringTok{"invest"}\NormalTok{, }\OperatorTok{**}\NormalTok{\{}\OperatorTok{**}\NormalTok{MODEL, }\OperatorTok{**}\NormalTok{infer, }\OperatorTok{**}\NormalTok{changes\})}


\BuiltInTok{print}\NormalTok{(full().gamma\_ci, full(level}\OperatorTok{=}\DecValTok{90}\NormalTok{).gamma\_ci)       }\CommentTok{\# 95 \% and 90 \% sets (hulls)}
\BuiltInTok{print}\NormalTok{(full(linearity\_test}\OperatorTok{=}\VariableTok{False}\NormalTok{).scalars[}\StringTok{"pval\_lin"}\NormalTok{])   }\CommentTok{\# nan: the test was skipped}
\BuiltInTok{print}\NormalTok{(full(continuity\_test}\OperatorTok{=}\VariableTok{True}\NormalTok{).scalars[}\StringTok{"pval\_cont"}\NormalTok{])  }\CommentTok{\# bootstrap test of continuity}

\CommentTok{\# Unit resampling (Algorithm 1 of Gong and Seo); slower, so fewer draws here.}
\NormalTok{unit }\OperatorTok{=}\NormalTok{ pyx.xtdpthresh(df, }\StringTok{"invest"}\NormalTok{, }\OperatorTok{**}\NormalTok{\{}\OperatorTok{**}\NormalTok{MODEL, }\StringTok{"n\_boot"}\NormalTok{: }\DecValTok{49}\NormalTok{, }\StringTok{"grid\_ci"}\NormalTok{: }\DecValTok{20}\NormalTok{, }\StringTok{"seed"}\NormalTok{: }\DecValTok{1}\NormalTok{,}
                                       \StringTok{"boot\_type"}\NormalTok{: }\StringTok{"unit"}\NormalTok{\})}
\BuiltInTok{print}\NormalTok{(unit.gamma\_ci, unit.scalars[}\StringTok{"ci\_nseg"}\NormalTok{])}

\CommentTok{\# Legacy coefficient bootstrap: a note is issued; the intervals are kept on the fit.}
\NormalTok{legacy }\OperatorTok{=}\NormalTok{ full(coef\_boot}\OperatorTok{=}\StringTok{"onestep"}\NormalTok{, coef\_ci\_type}\OperatorTok{=}\StringTok{"percentile"}\NormalTok{)}
\BuiltInTok{print}\NormalTok{(legacy.fit.bootstrap.coefficient.interval.shape)   }\CommentTok{\# (2, 9): lower and upper, e(b\_bootci)}
\end{Highlighting}
\end{Shaded}

\texttt{verbose} and \texttt{export\_gmm} are accepted and do nothing
(Section 11).

\begin{center}\rule{0.5\linewidth}{0.5pt}\end{center}

\section{4. Choosing your settings}\label{choosing-your-settings}

No setting is right in general: the data and the question decide. This
section says what each choice does according to the help file of the
command, and shows on the Hansen panel how far the answer can move. The
numbers are one panel and one model. They illustrate sensitivity; they
are not rules.

\subsection{Instruments: the choice that matters
most}\label{instruments-the-choice-that-matters-most}

Without \texttt{collapse} and \texttt{maxlag}, the number of instruments
grows with the square of the number of periods. Many instruments
relative to the number of units makes the Hansen test unreliable
(Roodman 2009) and, with the default centered covariance, tends to make
it over-reject (help file). \texttt{collapse} and \texttt{maxlag} shrink
the set, at the price of weaker identification of γ.

\begin{Shaded}
\begin{Highlighting}[]
\ImportTok{import}\NormalTok{ pandas }\ImportTok{as}\NormalTok{ pd}

\NormalTok{settings }\OperatorTok{=}\NormalTok{ \{}
    \StringTok{"collapsed, lags 3{-}8"}\NormalTok{: \{\},}
    \StringTok{"collapsed, lags 2{-}10"}\NormalTok{: \{}\StringTok{"maxlag"}\NormalTok{: (}\DecValTok{2}\NormalTok{, }\DecValTok{10}\NormalTok{)\},}
    \StringTok{"collapsed, lags 3{-}5"}\NormalTok{: \{}\StringTok{"maxlag"}\NormalTok{: (}\DecValTok{3}\NormalTok{, }\DecValTok{5}\NormalTok{)\},}
    \StringTok{"not collapsed, lags 3{-}5"}\NormalTok{: \{}\StringTok{"collapse"}\NormalTok{: }\VariableTok{False}\NormalTok{, }\StringTok{"maxlag"}\NormalTok{: (}\DecValTok{3}\NormalTok{, }\DecValTok{5}\NormalTok{)\},}
    \StringTok{"not collapsed, all lags"}\NormalTok{: \{}\StringTok{"collapse"}\NormalTok{: }\VariableTok{False}\NormalTok{, }\StringTok{"maxlag"}\NormalTok{: }\VariableTok{None}\NormalTok{\},}
\NormalTok{\}}
\NormalTok{rows }\OperatorTok{=}\NormalTok{ \{\}}
\ControlFlowTok{for}\NormalTok{ label, changes }\KeywordTok{in}\NormalTok{ settings.items():}
\NormalTok{    r }\OperatorTok{=}\NormalTok{ fit(}\OperatorTok{**}\NormalTok{changes)}
\NormalTok{    rows[label] }\OperatorTok{=}\NormalTok{ \{}\StringTok{"gamma"}\NormalTok{: r.gamma, }\StringTok{"instruments"}\NormalTok{: r.n\_instruments,}
                   \StringTok{"hansen\_p"}\NormalTok{: r.scalars[}\StringTok{"hansen\_p"}\NormalTok{], }\StringTok{"ar2\_p"}\NormalTok{: r.scalars[}\StringTok{"ar2\_p"}\NormalTok{]\}}
\BuiltInTok{print}\NormalTok{(pd.DataFrame(rows).T.}\BuiltInTok{round}\NormalTok{(}\DecValTok{4}\NormalTok{))}
\end{Highlighting}
\end{Shaded}

On this panel (560 units) the table reads:

{\def\LTcaptype{none} % do not increment counter
\begin{longtable}[]{@{}
  >{\raggedright\arraybackslash}p{(\linewidth - 8\tabcolsep) * \real{0.3810}}
  >{\raggedleft\arraybackslash}p{(\linewidth - 8\tabcolsep) * \real{0.1429}}
  >{\raggedleft\arraybackslash}p{(\linewidth - 8\tabcolsep) * \real{0.1905}}
  >{\raggedleft\arraybackslash}p{(\linewidth - 8\tabcolsep) * \real{0.1429}}
  >{\raggedleft\arraybackslash}p{(\linewidth - 8\tabcolsep) * \real{0.1429}}@{}}
\toprule\noalign{}
\begin{minipage}[b]{\linewidth}\raggedright
Setting
\end{minipage} & \begin{minipage}[b]{\linewidth}\raggedleft

\end{minipage} & \begin{minipage}[b]{\linewidth}\raggedleft
Instruments
\end{minipage} & \begin{minipage}[b]{\linewidth}\raggedleft
Hansen p
\end{minipage} & \begin{minipage}[b]{\linewidth}\raggedleft
AR(2) p
\end{minipage} \\
\midrule\noalign{}
\endhead
\bottomrule\noalign{}
\endlastfoot
collapsed, lags 3-8 & 0.4194 & 25 & 0.382 & 0.738 \\
collapsed, lags 2-10 & 0.0121 & 37 & 0.626 & 0.013 \\
collapsed, lags 3-5 & 0.5895 & 13 & 0.972 & 0.476 \\
not collapsed, lags 3-5 & 0.1216 & 144 & 0.000 & 0.939 \\
not collapsed, all lags & 0.1092 & 377 & 0.000 & 0.922 \\
\end{longtable}
}

The threshold estimate moves from 0.01 to 0.59 with the instrument set.
Two readings are grounded in the help file. The lag-2 window rejects
AR(2) (p = 0.013), which "casts doubt on the lag-2 instruments or on the
specification" (third example of the help file). The uncollapsed sets
have 144 and 377 instruments for 560 units and a Hansen p-value of
0.000, which the help file warns is unreliable in that situation: it is
not by itself evidence against the model, but the instability of  is a
fact to report. \texttt{pyx.robustness()} (Section 7) runs this kind of
variation for you.

A workable order: start from a collapsed set with a lag window justified
by the timing of the variables; check AR(2) and Hansen; move the window
and the collapsing; report what moves.

\subsection{First differences or forward orthogonal
deviations}\label{first-differences-or-forward-orthogonal-deviations}

\texttt{method="fd"} is the transformation of Seo and Shin (2016) and of
\texttt{xthenreg}; \texttt{method="fod"} removes the unit effect by
forward orthogonal deviations. With interior gaps, FOD discards fewer
equations (an interior gap removes three FD equations and two FOD
equations). Asymptotic theory exists for FD; for FOD the authors of the
command report Monte Carlo evidence only (Nguyen and Lai 2026). On this
balanced panel both give  = 0.4194 (AR(2) p-values 0.738 and 0.944). On
an unbalanced panel, estimate both and put them side by side with
\texttt{pyx.compare} (example 3 does exactly that); \texttt{diagnose()}
also points it out when FD is used on an unbalanced panel.

\subsection{Jump or kink}\label{jump-or-kink}

The default allows a discontinuity at γ. \texttt{kink=True} imposes
continuity; \texttt{q} should then be a regressor, or the model is a
one-sided hinge and a warning says so. If the function is continuous,
the unrestricted estimator of γ converges at rate n\^{}(1/4) to a
non-normal limit, whereas the kink estimator is asymptotically normal
(Gong and Seo 2026), so the choice affects inference, not only the fit.
\texttt{continuity\_test=True} tests the kink model against the jump
model. On the Hansen panel the kink model gives  = 0.032 and a Hansen
p-value of 0.000, against 0.4194 and 0.382 for the jump model.

\subsection{Threshold search}\label{threshold-search}

The defaults are a uniform grid of 100 points between the 5th and 95th
percentile of \texttt{q}. Uniform grids of 100, 199 and 397 points are
nested, so \texttt{grid=199} or \texttt{grid=397} is the natural check
that  does not depend on the grid. \texttt{grid\_type="quantile"} puts
points where the data are; \texttt{refine} adds the observed values of
\texttt{q} next to ; \texttt{trim} moves the ends; \texttt{min\_regime}
demands more observations in each regime. On the Hansen panel these
choices keep  between 0.414 and 0.421 (trim 0.20: 0.4144; quantile
grid: 0.4210; observed support: 0.4137; refine 3: 0.4150).

\subsection{Bootstrap}\label{bootstrap}

\texttt{n\_boot=299} and \texttt{grid\_ci=100} are the defaults. The
smallest possible p-value is 1/(\texttt{n\_boot} + 1), so a very small
\texttt{n\_boot} cannot show a small p-value. For exploring,
\texttt{n\_boot=99,\ grid\_ci=40} is much faster; use the defaults, or
more, for what you report. \texttt{boot\_type="unit"} follows the
algorithm with a proof (Gong and Seo 2026) but is much slower. Always
fix a \texttt{seed} for results you publish (a seed is recorded anyway).

\subsection{Variance}\label{variance}

\texttt{vce="robust"} and \texttt{vce="windmeijer"} give the same
estimates and thresholds; only the variance of the slopes and the AR
statistics that include  change. \texttt{center=False} reproduces
\texttt{xtabond2}\textquotesingle s uncentered covariance (and its
Hansen statistic).

\subsection{Before you report}\label{before-you-report}

\begin{enumerate}
\def\labelenumi{\arabic{enumi}.}
\tightlist
\item
  Keep \texttt{res.seed} and \texttt{res.settings\_log()}, or
  \texttt{res.to\_json()}.
\item
  Check the grid (\texttt{grid=199}) and one other instrument set.
\item
  Run \texttt{pyx.robustness()} and report how stable the threshold is.
\item
  Report the segments of the confidence set, not only its hull.
\item
  Report Hansen and AR(2) as diagnostics, with the number of instruments
  next to the number of units.
\end{enumerate}

\begin{center}\rule{0.5\linewidth}{0.5pt}\end{center}

\section{\texorpdfstring{5. Calling
\texttt{xthenreg}}{5. Calling xthenreg}}\label{calling-xthenreg}

\begin{Shaded}
\begin{Highlighting}[]
\NormalTok{res = pyx.xthenreg(data, y, x=(), *, q, endogenous=(), exogenous=(), inst=(),}
\NormalTok{                   id=None, time=None, where=None, static=False, kink=False,}
\NormalTok{                   grid\_num=100, trim\_rate=0.4, h\_0=1.5, boot\_reps=0, seed=None, draws=None,}
\NormalTok{                   engine="fast")}
\end{Highlighting}
\end{Shaded}

\texttt{xthenreg} implements the first-difference estimator of Seo, Kim
and Kim (2019) for strongly balanced panels, with the sup-Wald linearity
test of Seo and Shin (2016).

{\def\LTcaptype{none} % do not increment counter
\begin{longtable}[]{@{}
  >{\raggedright\arraybackslash}p{(\linewidth - 6\tabcolsep) * \real{0.1100}}
  >{\raggedright\arraybackslash}p{(\linewidth - 6\tabcolsep) * \real{0.1900}}
  >{\raggedright\arraybackslash}p{(\linewidth - 6\tabcolsep) * \real{0.0700}}
  >{\raggedright\arraybackslash}p{(\linewidth - 6\tabcolsep) * \real{0.6300}}@{}}
\toprule\noalign{}
\begin{minipage}[b]{\linewidth}\raggedright
Python
\end{minipage} & \begin{minipage}[b]{\linewidth}\raggedright
Stata
\end{minipage} & \begin{minipage}[b]{\linewidth}\raggedright
Default
\end{minipage} & \begin{minipage}[b]{\linewidth}\raggedright
Notes
\end{minipage} \\
\midrule\noalign{}
\endhead
\bottomrule\noalign{}
\endlastfoot
\texttt{y}, \texttt{x} & \texttt{depvar}, \texttt{indepvars} & & In
Stata the threshold variable is the \textbf{second word} of the variable
list, \texttt{xthenreg\ y\ q\ x1}; here it is the separate argument
\texttt{q}, so \texttt{xthenreg\ y\ q\ x1} is
\texttt{xthenreg(df,\ "y",\ {[}"x1"{]},\ q="q")}. \\
\texttt{endogenous} & \texttt{endogenous(varlist)} & none & Excluded
from \texttt{x}, as in Stata. \\
\texttt{exogenous} & \texttt{exogenous(varlist)} & none & Additional
exogenous regressors. Accepted by the command; not described in its help
file. \\
\texttt{inst} & \texttt{inst(varlist)} & none & Additional
instruments. \\
\texttt{static} & \texttt{static} & \texttt{False} & No automatic lag of
\texttt{y}. \\
\texttt{kink} & \texttt{kink} & \texttt{False} & Kink model. \\
\texttt{grid\_num} & \texttt{grid\_num(integer)} & \texttt{100} & Points
of the grid. \\
\texttt{trim\_rate} & \texttt{trim\_rate(real)} & \texttt{0.4} & Total
trimmed share of \texttt{q}. \\
\texttt{boot\_reps} & \texttt{boost(integer)} & \texttt{0} &
Replications of the bootstrap linearity test; \texttt{0} skips it. \\
\texttt{seed} & (set by \texttt{set\ seed}) & drawn & Seed of the
bootstrap weights; recorded in \texttt{res.seed}. \\
\texttt{draws} & (none) & \texttt{None} & Standard normal weights
\texttt{(boot\_reps,\ N)} to use instead of the generator. \\
\texttt{engine} & (none) & \texttt{"fast"} & \texttt{"fast"} or
\texttt{"literal"} moment computation; both give the same numbers,
\texttt{"literal"} follows the Mata code line by line. \\
\texttt{h\_0} & \texttt{h\_0(real)} & \texttt{1.5} & Multiplier of the
kernel bandwidth \texttt{h\ =\ h\_0\ ·\ 1.06\ ·\ N\^{}(-1/5)\ ·\ sd(q)}
of the covariance (jump model only; the kink model has no bandwidth).
Checked against Stata with \texttt{h\_0(3)} and \texttt{h\_0(0.7)}. \\
\end{longtable}
}

\begin{Shaded}
\begin{Highlighting}[]
\NormalTok{sk }\OperatorTok{=}\NormalTok{ pyx.xthenreg(df, }\StringTok{"invest"}\NormalTok{, [}\StringTok{"cashflow"}\NormalTok{, }\StringTok{"tobin\_q"}\NormalTok{], q}\OperatorTok{=}\StringTok{"debt"}\NormalTok{, }\BuiltInTok{id}\OperatorTok{=}\StringTok{"firm"}\NormalTok{, time}\OperatorTok{=}\StringTok{"year"}\NormalTok{,}
\NormalTok{                  boot\_reps}\OperatorTok{=}\DecValTok{49}\NormalTok{, seed}\OperatorTok{=}\DecValTok{3}\NormalTok{)}
\BuiltInTok{print}\NormalTok{(sk.names)}
\BuiltInTok{print}\NormalTok{(}\BuiltInTok{round}\NormalTok{(sk.gamma, }\DecValTok{4}\NormalTok{), }\BuiltInTok{round}\NormalTok{(sk.scalars[}\StringTok{"sup\_wald"}\NormalTok{], }\DecValTok{2}\NormalTok{), sk.scalars[}\StringTok{"boots\_p"}\NormalTok{])}
\BuiltInTok{print}\NormalTok{(sk.ci.}\BuiltInTok{round}\NormalTok{(}\DecValTok{3}\NormalTok{))          }\CommentTok{\# e(CI): lower and upper limits of every coefficient, γ last}
\BuiltInTok{print}\NormalTok{(pyx.compare([sk, full()], model\_names}\OperatorTok{=}\NormalTok{[}\StringTok{"xthenreg"}\NormalTok{, }\StringTok{"xtdpthresh"}\NormalTok{]).to\_text())}
\end{Highlighting}
\end{Shaded}

\textbf{Coefficient names} are those of the command. \texttt{Lag\_y\_b}
is the slope of \texttt{L.y} below the threshold and \texttt{x\_b} that
of a regressor; \texttt{cons\_d}, \texttt{Lag\_y\_d} and \texttt{x\_d}
are the changes above it (\texttt{\_d} for δ; \texttt{cons\_d} is the
intercept shift); \texttt{r} is the threshold. With \texttt{kink=True}
the names are \texttt{Lag\_y\_b}, \texttt{x\_b}, \texttt{kink\_slope}
and \texttt{r}; with \texttt{static=True} there is no
\texttt{Lag\_y\_*}. The table of the result shows them in the same three
blocks (lower, change, upper) as \texttt{xtdpthresh}.

\textbf{Inference.} \texttt{xthenreg} reports an asymptotic interval for
every coefficient, the threshold included (\texttt{res.ci},
\texttt{e(CI)}), not a bootstrap confidence set. Its linearity test is
the sup-Wald statistic of Seo and Shin (2016)
(\texttt{scalars{[}"sup\_wald"{]}}) with a bootstrap p-value
\texttt{boots\_p} that is \texttt{-1} when the test was not run
(\texttt{boot\_reps=0}) and can be exactly \texttt{0} when no
replication reaches the sample statistic. \texttt{scalars} also holds
\texttt{N} (units), \texttt{T}, \texttt{grid}, \texttt{trim},
\texttt{bs} (replications), \texttt{gamma}, \texttt{L} (number of moment
conditions) and \texttt{bandwidth} (kernel bandwidth of the covariance).

\textbf{What the result has.} Estimates, \texttt{se}, \texttt{cov},
\texttt{ci}, \texttt{gamma}, \texttt{seed}, \texttt{scalars},
\texttt{summary()}, and all the table and export methods
(\texttt{to\_latex}, \texttt{to\_html}, \ldots, \texttt{compare}).
Graphs, \texttt{diagnose()}, \texttt{settings\_log()}, \texttt{predict},
\texttt{augment} and \texttt{ci\_table()} belong to \texttt{xtdpthresh}
results only; \texttt{pyx.plot} called on an \texttt{xthenreg} result
raises a \texttt{SpecificationError} (key \texttt{plot\_xthenreg}) that
says so.

\textbf{Against \texttt{xtdpthresh}.} Same moment conditions and weight
matrices under \texttt{method="fd"}, so the same estimates on the same
grid; but the default trimming is 0.4 instead of 0.10, the threshold
variable is given differently (see the table above), \texttt{xthenreg}
needs a strongly balanced panel and has neither \texttt{fod},
\texttt{td}, \texttt{maxlag}, \texttt{collapse} nor \texttt{iv}, and
strictly exogenous regressors are instrumented by their levels where
\texttt{xtdpthresh} uses their transformed values.

\begin{center}\rule{0.5\linewidth}{0.5pt}\end{center}

\section{6. Reading the output}\label{reading-the-output}

\subsection{The printed result}\label{the-printed-result}

\texttt{print(res)} (or just \texttt{res} in a notebook) shows the
sample, the threshold with its confidence interval, the three blocks of
slopes, and the four diagnostic tests:

\begin{verbatim}
Dynamic panel threshold model
Dependent variable: invest   Units: 560   Observations: 6720
Threshold variable: debt   Threshold  = 0.4194   Threshold 95 % interval: [0.0122, 0.5227]
Threshold effect: jump   Transformation: FD   Estimator: two-step
...
Lower regime      slopes below the threshold                       (e(b), equation "lower")
Change at the threshold   shifts above it: _cons is δ₁, the rest δ₂  (e(b), equation "change")
Upper regime      lower + change, with its own standard error      (e(b_upper), e(V_upper))
...
Hansen J (p-value): 0.3818   AR(1) (p-value): 0.0000
AR(2) (p-value): 0.7382   Linearity test (p-value): 0.0067
\end{verbatim}

"Observations" is the number of stacked equations (\texttt{e(N)}), not
of firm-years. Under \texttt{kink} the change block holds only the
change in the slope of \texttt{q}, and the upper block only the slope of
\texttt{q} above the threshold (\texttt{res.upper} is \texttt{None} when
\texttt{q} is not a regressor).

\subsection{Numbers as objects}\label{numbers-as-objects}

\begin{Shaded}
\begin{Highlighting}[]
\BuiltInTok{print}\NormalTok{(res.params)          }\CommentTok{\# Series indexed \textquotesingle{}lower:cashflow\textquotesingle{}, \textquotesingle{}change:\_cons\textquotesingle{}, ... (e(b))}
\BuiltInTok{print}\NormalTok{(res.se)              }\CommentTok{\# standard errors from the reported variance (sqrt of the diagonal of e(V))}
\BuiltInTok{print}\NormalTok{(res.cov.shape)       }\CommentTok{\# e(V)}
\BuiltInTok{print}\NormalTok{(res.upper)           }\CommentTok{\# upper{-}regime slopes and standard errors (e(b\_upper), e(V\_upper))}
\BuiltInTok{print}\NormalTok{(res.gamma, res.gamma\_ci)}
\BuiltInTok{print}\NormalTok{(res.to\_dataframe())  }\CommentTok{\# estimate, se, z, p and the interval, indexed by (block, variable)}
\end{Highlighting}
\end{Shaded}

\texttt{res.gamma\_ci} is the \textbf{hull} of the confidence set: the
smallest and largest accepted candidate. When the set has several
segments (\texttt{res.scalars{[}"ci\_nseg"{]}\ \textgreater{}\ 1}), the
range between them contains rejected values, so look at
\texttt{res.ci\_table()} or \texttt{res.plot("ci")} and report the
segments.

\subsection{\texorpdfstring{\texttt{res.scalars}}{res.scalars}}\label{res.scalars}

A dictionary with the names of Stata\textquotesingle s \texttt{e()}
(\texttt{nan} when a value does not exist for the fit):

{\def\LTcaptype{none} % do not increment counter
\begin{longtable}[]{@{}
  >{\raggedright\arraybackslash}p{(\linewidth - 2\tabcolsep) * \real{0.2000}}
  >{\raggedright\arraybackslash}p{(\linewidth - 2\tabcolsep) * \real{0.8000}}@{}}
\toprule\noalign{}
\begin{minipage}[b]{\linewidth}\raggedright
Group
\end{minipage} & \begin{minipage}[b]{\linewidth}\raggedright
Names
\end{minipage} \\
\midrule\noalign{}
\endhead
\bottomrule\noalign{}
\endlastfoot
Sample & \texttt{N} (equations), \texttt{N\_units}, \texttt{N\_clust},
\texttt{N\_switch} (units whose regime changes within their equations at
) \\
Instruments & \texttt{N\_iv}, \texttt{N\_iv\_common},
\texttt{N\_iv\_dep}, \texttt{N\_iv\_dep\_near}, \texttt{iv\_dep\_res}
(columns dropped as common to all units or numerically dependent, and
the largest residual of a dropped one) \\
Estimator & \texttt{gamma}, \texttt{obj} (criterion),
\texttt{estimator\_twostep} (0 = one-step fallback),
\texttt{W1\_fallback}, \texttt{level}, \texttt{joint\_vce},
\texttt{vce\_applied}, \texttt{wind\_same\_threshold},
\texttt{gamma\_bw}, \texttt{q\_nvals\_bw}, \texttt{kink\_refined} \\
Grid & \texttt{q\_lo}, \texttt{q\_hi} (trimming bounds),
\texttt{grid\_admitted} \\
Hansen & \texttt{hansen}, \texttt{hansen\_df}, \texttt{hansen\_p} \\
Serial correlation & \texttt{ar1}, \texttt{ar1\_p}, \texttt{ar2},
\texttt{ar2\_p}, \texttt{ar\_joint}, \texttt{ar1\_cond},
\texttt{ar2\_cond}; \texttt{ar1\_b0}, \texttt{ar1\_T1},
\texttt{ar1\_TT}, \texttt{ar2\_b0}, \texttt{ar2\_T1}, \texttt{ar2\_TT}
are the pieces of the statistics, kept for checking against Stata \\
Confidence set & \texttt{gamma\_lo}, \texttt{gamma\_hi},
\texttt{ci\_nseg}, \texttt{ci\_empty}, \texttt{ci\_incomplete},
\texttt{ci\_unresolved}, \texttt{boundary\_warn} \\
Tests & \texttt{pval\_lin}, \texttt{boot\_linearity\_valid},
\texttt{pval\_cont}, \texttt{boot\_continuity\_valid},
\texttt{citest\_gamma}, \texttt{citest\_D}, \texttt{citest\_crit},
\texttt{citest\_accept}, \texttt{citest\_p}, \texttt{citest\_status},
\texttt{citest\_draws}, \texttt{rseed} \\
\end{longtable}
}

The tables behind the confidence set and the regimes:

\begin{Shaded}
\begin{Highlighting}[]
\BuiltInTok{print}\NormalTok{(res.ci\_table().head())   }\CommentTok{\# gamma, statistic, critical, accepted, n\_valid, status}
\BuiltInTok{print}\NormalTok{(res.regime\_table())      }\CommentTok{\# observations, units, share and the mean of each variable by regime}
\end{Highlighting}
\end{Shaded}

In \texttt{ci\_table()} the \texttt{status} is 1 for a resolved test, 2
for accepted with a zero statistic, 3 for a candidate that is not
admissible, and 4 to 6 for unresolved ones. Statuses 3 to 6 are not
rejections.

\subsection{A reading order that
works}\label{a-reading-order-that-works}

\begin{enumerate}
\def\labelenumi{\arabic{enumi}.}
\tightlist
\item
  \textbf{The sample.} \texttt{n\_units}, \texttt{n\_obs} and
  \texttt{n\_instruments} (or \texttt{res.settings\_log()}): were units
  dropped, and is the number of instruments small next to the number of
  units? Many instruments relative to units makes the Hansen test
  unreliable (Roodman 2009).
\item
  \textbf{The threshold.}  is a grid point (a point between grid points
  under \texttt{kink} or \texttt{refine}). It has no standard error;
  inference rests on the confidence set. Check the number of segments,
  \texttt{boundary\_warn} (a bound at the edge of the confidence-set
  grid: 1 lower, 2 upper, 3 both), \texttt{ci\_incomplete} (no interval
  is reported if some candidates are unresolved) and \texttt{ci\_empty}.
\item
  \textbf{Linearity.} The test of H₀: δ = 0. A rejection says a
  threshold effect exists; it does not say where the threshold is. The
  p-value is (1 + R)/(1 + B).
\item
  \textbf{Hansen J.} A diagnostic, with \texttt{N\_iv\ -\ k\ -\ 1}
  degrees of freedom. Its chi-square reference presumes regular
  identification, which fails near a continuous threshold. A large
  p-value does not show that the instruments are valid. With
  \texttt{center=False} it equals the statistic of \texttt{xtabond2} for
  the same instruments.
\item
  \textbf{AR(1) and AR(2).} AR(1) is expected to reject. A rejection of
  AR(2) indicates serial correlation in the errors, which invalidates
  the lagged instruments: reconsider \texttt{maxlag} or the
  specification (the second example of the Stata help file shows this
  with \texttt{maxlag(2\ 10)}).
\item
  \textbf{The fallbacks.} \texttt{estimator\_twostep\ ==\ 0},
  \texttt{W1\_fallback}, \texttt{vce\_applied\ ==\ 0} mean the command
  fell back to a simpler computation; \texttt{diagnose()} and the
  printed notes say which.
\item
  \textbf{Discrete \texttt{q}.} In the jump model, the joint variance
  and the AR statistics with  assume that \texttt{q} has a continuous
  density at the threshold. \texttt{q\_nvals\_bw} below 10 triggers a
  warning; the conditional statistics \texttt{ar1\_cond} and
  \texttt{ar2\_cond} do not use that assumption.
\end{enumerate}

\subsection{\texorpdfstring{\texttt{diagnose()} and
\texttt{settings\_log()}}{diagnose() and settings\_log()}}\label{diagnose-and-settings_log}

\texttt{res.diagnose()} lists, in plain language, what deserves
attention, each finding with a recommendation and a level (warning or
notice). It checks: the number of instruments against the number of
units; a Hansen test that rejects or sits suspiciously close to one, and
an AR(2) test that rejects; a threshold at an edge of the search grid;
the confidence set (empty, incomplete, touching the edge of its grid,
several segments, very wide, or so narrow that the grid is too coarse);
a flat criterion profile or a distant rival threshold that fits almost
as well; a small regime or a threshold variable with few distinct
values; first differences on an unbalanced panel; the one-step and
first-weight fallbacks; and the result of the continuity test when it
was run. In English or French: \texttt{res.diagnose(lang="fr")}.

\texttt{res.settings\_log()} is the journal of what the estimation
\emph{really} used: the sample (units with equations, dropped units),
the model, the instruments kept and removed, the first-step weight, the
grid (requested, evaluated, admitted, range, smallest regime allowed),
the variance estimator, the bootstrap (type, draws, seed, points of the
grid) and the package version. Keep it with the results.

\begin{Shaded}
\begin{Highlighting}[]
\BuiltInTok{print}\NormalTok{(res.diagnose())}
\BuiltInTok{print}\NormalTok{(res.settings\_log())}
\end{Highlighting}
\end{Shaded}

\subsection{Two diagnoses, read}\label{two-diagnoses-read}

\emph{Serial correlation.} The same model with instruments from lag 2
(the third example of the Stata help file):

\begin{Shaded}
\begin{Highlighting}[]
\NormalTok{ar2 }\OperatorTok{=}\NormalTok{ full(maxlag}\OperatorTok{=}\NormalTok{(}\DecValTok{2}\NormalTok{, }\DecValTok{10}\NormalTok{))}
\BuiltInTok{print}\NormalTok{(ar2.diagnose())}
\end{Highlighting}
\end{Shaded}

\begin{Shaded}
\begin{Highlighting}[]
\NormalTok{Diagnostic report: 1 warning(s), 1 notice(s)}

\NormalTok{[warning] The AR(2) test rejects the absence of second{-}order serial correlation in the differenced residuals (p = 0.013): lagged levels may not be valid instruments.}
\NormalTok{    Recommendation: Start the instruments one lag deeper (for example maxlag=(3, 8)) or add dynamics.}
\NormalTok{[notice] The confidence set of the threshold reaches the lower edge of the candidates: it is cut by the grid, not by the data.}
\NormalTok{    Recommendation: Widen the grid (lower trim) before reading the interval.}
\end{Highlighting}
\end{Shaded}

The warning is the one to act on: the instruments are suspect, so the
estimate ( = 0.0121, near the lower end of the grid) should not be
reported as it stands. The notice says the confidence set touches the
edge of the candidates: it is cut by the grid, so do not read its end as
a bound.

\emph{Too many instruments.} The same model on the first 79 firms,
without \texttt{collapse} and with every lag:

\begin{Shaded}
\begin{Highlighting}[]
\NormalTok{small }\OperatorTok{=}\NormalTok{ pyx.xtdpthresh(df[df[}\StringTok{"firm"}\NormalTok{] }\OperatorTok{\textless{}} \DecValTok{80}\NormalTok{], }\StringTok{"invest"}\NormalTok{,}
                       \OperatorTok{**}\NormalTok{\{}\OperatorTok{**}\NormalTok{MODEL, }\OperatorTok{**}\NormalTok{infer, }\StringTok{"collapse"}\NormalTok{: }\VariableTok{False}\NormalTok{, }\StringTok{"maxlag"}\NormalTok{: }\VariableTok{None}\NormalTok{\})}
\BuiltInTok{print}\NormalTok{(small.n\_units, small.n\_instruments, small.scalars[}\StringTok{"estimator\_twostep"}\NormalTok{])}
\BuiltInTok{print}\NormalTok{(small.diagnose())}
\end{Highlighting}
\end{Shaded}

\begin{Shaded}
\begin{Highlighting}[]
\NormalTok{79 377 0.0}
\NormalTok{Diagnostic report: 2 warning(s), 3 notice(s)}

\NormalTok{[warning] 377 instruments for 79 units: with more instruments than units the Hansen test loses its power and the weight matrix is poorly estimated.}
\NormalTok{    Recommendation: Reduce the instruments: collapse=True or limit the lags with maxlag=(2, 4).}
\NormalTok{[warning] The estimate stays at one step: the second{-}step weight could not be used.}
\NormalTok{    Recommendation: Reduce the instruments (collapse=True or limit the lags with maxlag=(2, 4)); the Hansen test is not available at one step.}
\NormalTok{[notice] The confidence set of the threshold reaches the lower and upper edge of the candidates: it is cut by the grid, not by the data.}
\NormalTok{    Recommendation: Widen the grid (lower trim) before reading the interval.}
\NormalTok{[notice] The confidence set covers 100 \% of the search grid: the threshold is weakly identified.}
\NormalTok{    Recommendation: Treat the threshold with caution; more units, a longer panel or a more variable debt would help.}
\NormalTok{[notice] The criterion profile is flat: 100 \% of the thresholds are within 3.8 of the minimum.}
\NormalTok{    Recommendation: The estimated threshold is not sharply determined; compare it with the confidence set and plot(fit, \textquotesingle{}profile\textquotesingle{}).}
\end{Highlighting}
\end{Shaded}

Here the cause is upstream of everything else: with more instruments
than units the second-step weight could not be used, the command falls
back to one step (\texttt{estimator\_twostep} is 0, and there is no
Hansen test), and the confidence set is the whole grid. The five
findings share one remedy, fewer instruments.

\begin{center}\rule{0.5\linewidth}{0.5pt}\end{center}

\section{7. Working with a result}\label{working-with-a-result}

\textbf{Predictions and residuals.} \texttt{res.predict(kind,\ frame)}
returns a Series for the estimation data; \texttt{kind} is
\texttt{"residuals"}, \texttt{"xb"}, \texttt{"ar\_residuals"} (the
first-difference residuals used by the AR tests; they equal
\texttt{"residuals"} under FD) or \texttt{"regime"} (1 when
\texttt{q\ \textgreater{}\ }, missing where \texttt{q} is). They are
the series computed during the estimation, not recomputed. They are the
counterparts of Stata\textquotesingle s
\texttt{predict,\ residuals\ \textbar{}\ arresiduals\ \textbar{}\ xb\ \textbar{}\ regime}.

\texttt{res.augment(frame,\ kinds=…,\ prefix="",\ overwrite=False)}
returns a copy of the frame with those series as columns (default
\texttt{kinds}: residuals, xb, regime), named after the series with
\texttt{prefix} in front. A column that already has the name is an error
(\texttt{augment\_exists}) unless \texttt{overwrite=True}.

\begin{Shaded}
\begin{Highlighting}[]
\NormalTok{aug }\OperatorTok{=}\NormalTok{ res.augment(df, kinds}\OperatorTok{=}\NormalTok{(}\StringTok{"residuals"}\NormalTok{, }\StringTok{"regime"}\NormalTok{))}
\BuiltInTok{print}\NormalTok{(aug[}\StringTok{"regime"}\NormalTok{].mean())}
\NormalTok{both }\OperatorTok{=}\NormalTok{ res.augment(df, prefix}\OperatorTok{=}\StringTok{"m1\_"}\NormalTok{, overwrite}\OperatorTok{=}\VariableTok{True}\NormalTok{)}
\BuiltInTok{print}\NormalTok{([c }\ControlFlowTok{for}\NormalTok{ c }\KeywordTok{in}\NormalTok{ both.columns }\ControlFlowTok{if}\NormalTok{ c.startswith(}\StringTok{"m1\_"}\NormalTok{)])}
\end{Highlighting}
\end{Shaded}

\textbf{Tables.} \texttt{res.to\_latex(path)}, \texttt{to\_html},
\texttt{to\_markdown}, \texttt{to\_text}, \texttt{to\_csv},
\texttt{to\_word} and \texttt{to\_excel} (the last two need
\texttt{pip\ install\ "pyxtdpthresh{[}export{]}"}) write the academic
table and return its text. \texttt{res.table(**options)} returns the
table object, and
\texttt{pyx.compare({[}a,\ b,\ …{]},\ model\_names={[}…{]})} puts
several models side by side (\texttt{xthenreg} against
\texttt{xtdpthresh}, FD against FOD, jump against kink). The options,
shared by all of them:

{\def\LTcaptype{none} % do not increment counter
\begin{longtable}[]{@{}
  >{\raggedright\arraybackslash}p{(\linewidth - 2\tabcolsep) * \real{0.3800}}
  >{\raggedright\arraybackslash}p{(\linewidth - 2\tabcolsep) * \real{0.6200}}@{}}
\toprule\noalign{}
\begin{minipage}[b]{\linewidth}\raggedright
Option
\end{minipage} & \begin{minipage}[b]{\linewidth}\raggedright
Effect
\end{minipage} \\
\midrule\noalign{}
\endhead
\bottomrule\noalign{}
\endlastfoot
\texttt{statistic} & under each estimate: \texttt{"se"} (default),
\texttt{"t"}, \texttt{"p"} or \texttt{"ci"} (interval at
\texttt{level}) \\
\texttt{digits}, \texttt{star\_levels}, \texttt{level} & decimals (3),
significance levels (0.1, 0.05, 0.01), level of the intervals \\
\texttt{model\_names}, \texttt{column\_labels} & column names, and a
group label above them \\
\texttt{covariate\_labels}, \texttt{dep\_var\_label} & new names for
variables (\texttt{\{"cashflow":\ "Cash\ flow"\}}; a
\texttt{"block:variable"} key names one regime) and for the dependent
variable \\
\texttt{keep}, \texttt{omit}, \texttt{order} & regular expressions
selecting, dropping and ordering rows (against the name or
\texttt{block:variable}) \\
\texttt{add\_lines}, \texttt{notes} & extra lines under the coefficients
(\texttt{{[}("Firm\ effects",\ {[}"Yes",\ "Yes"{]}){]}}) and notes \\
\texttt{title}, \texttt{label}, \texttt{font\_size},
\texttt{single\_row} & caption, LaTeX label, LaTeX size command, one
line per estimate \\
\texttt{omit\_stat} & diagnostic lines to drop (\texttt{"n\_obs"},
\texttt{"hansen\_p"}, \ldots) \\
\texttt{lang} & \texttt{"en"} or \texttt{"fr"} for labels and the
decimal sign \\
\texttt{full\_output=True} & (of \texttt{res.table} and
\texttt{res.to\_*}) one more line per scalar of \texttt{e()} \\
\end{longtable}
}

\begin{Shaded}
\begin{Highlighting}[]
\NormalTok{fod }\OperatorTok{=}\NormalTok{ pyx.xtdpthresh(df, }\StringTok{"invest"}\NormalTok{, method}\OperatorTok{=}\StringTok{"fod"}\NormalTok{, bootstrap}\OperatorTok{=}\VariableTok{False}\NormalTok{, }\OperatorTok{**}\NormalTok{MODEL)}
\NormalTok{kink }\OperatorTok{=}\NormalTok{ pyx.xtdpthresh(df, }\StringTok{"invest"}\NormalTok{, kink}\OperatorTok{=}\VariableTok{True}\NormalTok{, bootstrap}\OperatorTok{=}\VariableTok{False}\NormalTok{, }\OperatorTok{**}\NormalTok{MODEL)}
\NormalTok{table }\OperatorTok{=}\NormalTok{ pyx.compare([fast, fod, kink], model\_names}\OperatorTok{=}\NormalTok{[}\StringTok{"FD jump"}\NormalTok{, }\StringTok{"FOD jump"}\NormalTok{, }\StringTok{"FD kink"}\NormalTok{])}
\BuiltInTok{print}\NormalTok{(table.to\_text())}

\NormalTok{custom }\OperatorTok{=}\NormalTok{ res.table(statistic}\OperatorTok{=}\StringTok{"ci"}\NormalTok{, digits}\OperatorTok{=}\DecValTok{2}\NormalTok{, keep}\OperatorTok{=}\NormalTok{[}\StringTok{"cashflow"}\NormalTok{, }\StringTok{"debt"}\NormalTok{],}
\NormalTok{                   covariate\_labels}\OperatorTok{=}\NormalTok{\{}\StringTok{"cashflow"}\NormalTok{: }\StringTok{"Cash flow"}\NormalTok{\},}
\NormalTok{                   add\_lines}\OperatorTok{=}\NormalTok{[(}\StringTok{"Firm effects"}\NormalTok{, [}\StringTok{"Yes"}\NormalTok{])], title}\OperatorTok{=}\StringTok{"Investment"}\NormalTok{)}
\BuiltInTok{print}\NormalTok{(custom.to\_text())}
\NormalTok{res.to\_latex(}\StringTok{"table.tex"}\NormalTok{, full\_output}\OperatorTok{=}\VariableTok{True}\NormalTok{)}
\end{Highlighting}
\end{Shaded}

\textbf{Plots.}
\texttt{res.plot(kind,\ path=None,\ lang=None,\ **options)} returns a
matplotlib \texttt{Figure} and writes it to \texttt{path}
(\texttt{.png}, \texttt{.pdf} or \texttt{.svg}) when given. They are
\texttt{xtdpthresh} results\textquotesingle{} methods (Section 5).

{\def\LTcaptype{none} % do not increment counter
\begin{longtable}[]{@{}
  >{\raggedright\arraybackslash}p{(\linewidth - 4\tabcolsep) * \real{0.1000}}
  >{\raggedright\arraybackslash}p{(\linewidth - 4\tabcolsep) * \real{0.5500}}
  >{\raggedright\arraybackslash}p{(\linewidth - 4\tabcolsep) * \real{0.3500}}@{}}
\toprule\noalign{}
\begin{minipage}[b]{\linewidth}\raggedright
\texttt{kind}
\end{minipage} & \begin{minipage}[b]{\linewidth}\raggedright
Draws
\end{minipage} & \begin{minipage}[b]{\linewidth}\raggedright
Options
\end{minipage} \\
\midrule\noalign{}
\endhead
\bottomrule\noalign{}
\endlastfoot
\texttt{profile} & the GMM criterion over the grid of thresholds & \\
\texttt{ci} & the grid-bootstrap statistic against its critical value,
accepted set shaded & \\
\texttt{linearity} & the bootstrap distribution of the linearity test
& \\
\texttt{coef} & slopes of both regimes and their change, with intervals
& \texttt{variables}, \texttt{free\_scale} \\
\texttt{regimes} & the distribution of \texttt{q} and the share in each
regime & \texttt{bins} (30), \texttt{trim} (0.01 cut from each tail) \\
\texttt{effect} & the effect of one regressor according to the regime &
\texttt{var}, \texttt{trim} \\
\texttt{switch} & the share of the upper regime by period, and the units
that change regime & \\
\texttt{instruments} & instrument columns by period against the number
of units & \\
\texttt{residuals} & transformed residuals by period and the AR(1) and
AR(2) tests & \\
\end{longtable}
}

\texttt{ci} and \texttt{linearity} need the bootstrap.
\texttt{robustness} takes a \texttt{RobustnessResult};
\texttt{coverage}, \texttt{power} and \texttt{distribution} take Monte
Carlo results (one, a list, or a \texttt{\{name:\ result\}} mapping)
through \texttt{pyx.plot(result,\ kind)}. Figures follow the style of
the article and come in English or French (\texttt{lang="fr"}).

\begin{Shaded}
\begin{Highlighting}[]
\NormalTok{fig }\OperatorTok{=}\NormalTok{ res.plot(}\StringTok{"profile"}\NormalTok{, path}\OperatorTok{=}\StringTok{"profile.png"}\NormalTok{)}
\NormalTok{res.plot(}\StringTok{"effect"}\NormalTok{, var}\OperatorTok{=}\StringTok{"cashflow"}\NormalTok{, path}\OperatorTok{=}\StringTok{"effect.png"}\NormalTok{)}
\NormalTok{res.plot(}\StringTok{"regimes"}\NormalTok{, bins}\OperatorTok{=}\DecValTok{20}\NormalTok{, lang}\OperatorTok{=}\StringTok{"fr"}\NormalTok{, path}\OperatorTok{=}\StringTok{"regimes\_fr.png"}\NormalTok{)}
\end{Highlighting}
\end{Shaded}

\textbf{Robustness.}
\texttt{pyx.robustness(df,\ y,\ x,\ q=…,\ settings=…,\ product=False,\ **options)}
re-estimates the model on a set of settings and gathers the threshold
with its confidence set, the slopes, the linearity p-value and the
diagnostics of each run. \texttt{settings} maps an option to the values
to try (default: \texttt{grid} 50/100/200, \texttt{trim} 0.05/0.10/0.15,
\texttt{collapse}, \texttt{method}, \texttt{kink}); the options of the
call are the reference. By default one option changes at a time;
\texttt{product=True} runs every combination. All runs share one
bootstrap seed, so the settings, not the draws, explain the differences.

\begin{Shaded}
\begin{Highlighting}[]
\NormalTok{rob }\OperatorTok{=}\NormalTok{ pyx.robustness(df, }\StringTok{"invest"}\NormalTok{, n\_boot}\OperatorTok{=}\DecValTok{99}\NormalTok{, grid\_ci}\OperatorTok{=}\DecValTok{40}\NormalTok{, seed}\OperatorTok{=}\DecValTok{1}\NormalTok{, }\OperatorTok{**}\NormalTok{MODEL)}
\BuiltInTok{print}\NormalTok{(rob.table[[}\StringTok{"gamma"}\NormalTok{, }\StringTok{"gamma\_lo"}\NormalTok{, }\StringTok{"gamma\_hi"}\NormalTok{, }\StringTok{"ci\_nseg"}\NormalTok{, }\StringTok{"in\_reference\_ci"}\NormalTok{]])}
\BuiltInTok{print}\NormalTok{(rob.stability())     }\CommentTok{\# share of the other runs whose threshold stays in the reference set}

\NormalTok{own }\OperatorTok{=}\NormalTok{ pyx.robustness(df, }\StringTok{"invest"}\NormalTok{, n\_boot}\OperatorTok{=}\DecValTok{49}\NormalTok{, grid\_ci}\OperatorTok{=}\DecValTok{20}\NormalTok{, seed}\OperatorTok{=}\DecValTok{1}\NormalTok{,}
\NormalTok{                     settings}\OperatorTok{=}\NormalTok{\{}\StringTok{"trim"}\NormalTok{: [}\FloatTok{0.05}\NormalTok{, }\FloatTok{0.15}\NormalTok{], }\StringTok{"method"}\NormalTok{: [}\StringTok{"fod"}\NormalTok{]\}, product}\OperatorTok{=}\VariableTok{True}\NormalTok{, }\OperatorTok{**}\NormalTok{MODEL)}
\BuiltInTok{print}\NormalTok{(own.table.index.tolist())    }\CommentTok{\# reference, then the combinations}
\NormalTok{rob.plot(path}\OperatorTok{=}\StringTok{"robustness.png"}\NormalTok{)}
\BuiltInTok{print}\NormalTok{(rob.coefficients().head())   }\CommentTok{\# every slope of every run, with its standard error}
\end{Highlighting}
\end{Shaded}

\texttt{rob.table} has one row per run;
\texttt{stability(use="segments")} counts a run as stable when its
threshold lies in an accepted segment of the reference set
(\texttt{use="hull"} uses the hull). A run that fails is kept in the
table with its message and does not stop the analysis; only a failure of
the reference does. \texttt{rob.stargazer()} returns the table of all
the runs side by side.

\textbf{Records.} \texttt{res.to\_json(path)} writes a reproducible
record of the fit with seven keys: \texttt{model},
\texttt{specification} (variables), \texttt{options} (every option, as
used), \texttt{scalars}, \texttt{ci\_grid} (the inversion table),
\texttt{coefficients} and \texttt{environment} (versions of the package,
Python, numpy, scipy, pandas and the platform). \texttt{res.fit} holds
everything the estimator computed, as frozen containers.

\begin{Shaded}
\begin{Highlighting}[]
\ImportTok{import}\NormalTok{ json}

\NormalTok{record }\OperatorTok{=}\NormalTok{ json.loads(res.to\_json())}
\BuiltInTok{print}\NormalTok{(}\BuiltInTok{list}\NormalTok{(record), record[}\StringTok{"options"}\NormalTok{][}\StringTok{"seed"}\NormalTok{])}
\end{Highlighting}
\end{Shaded}

\begin{center}\rule{0.5\linewidth}{0.5pt}\end{center}

\section{8. Reproducibility, speed and
memory}\label{reproducibility-speed-and-memory}

\textbf{Seeds.} Everything random in \texttt{xtdpthresh} comes from
\texttt{seed}. The Stata generator is reproduced (\texttt{set\ seed},
\texttt{uniform()}, \texttt{invnormal()}), so a seeded bootstrap draws
the same weights as Stata. A seed that is omitted is drawn and recorded
in \texttt{res.seed}, \texttt{res.settings\_log()} and
\texttt{res.to\_json()}.

\textbf{Global options.}

\begin{Shaded}
\begin{Highlighting}[]
\CommentTok{\# not run by the tests}
\NormalTok{pyx.set\_options(lang}\OperatorTok{=}\StringTok{"fr"}\NormalTok{)        }\CommentTok{\# messages, reports, figures in French}
\NormalTok{pyx.set\_options(threads}\OperatorTok{=}\DecValTok{4}\NormalTok{)        }\CommentTok{\# threads of the matrix products (default: half the logical cores)}
\NormalTok{pyx.set\_options(progress}\OperatorTok{=}\VariableTok{True}\NormalTok{)    }\CommentTok{\# progress line during long bootstraps, on standard error}
\ControlFlowTok{with}\NormalTok{ pyx.option\_context(lang}\OperatorTok{=}\StringTok{"fr"}\NormalTok{):}
\NormalTok{    ...                           }\CommentTok{\# temporary change}
\end{Highlighting}
\end{Shaded}

\texttt{set\_options} returns the previous options. Numba kernels run on
one thread, which is faster than letting them compete with the matrix
products.

\textbf{First call.} The Numba kernels are compiled the first time they
are used and cached on disk. \texttt{pyx.warmup()} pays that cost ahead
of time (a few seconds, once per environment).

\textbf{Where the time goes.} The bootstrap is 80--92 \% of a full run.
For a quick look at a specification, estimate with
\texttt{bootstrap=False}; increase \texttt{n\_boot} and
\texttt{grid\_ci} for the final run. With many instruments the cost is
in one matrix product that already runs close to the speed of the
machine; \texttt{collapse} and \texttt{maxlag} shrink it.

\textbf{Memory.} Memory grows with the number of equations times the
number of raw instrument columns. A note is issued before a fit expected
to need more than 2 GB, and a warning above 8 GB; \texttt{collapse},
\texttt{maxlag} and fewer draws reduce it.

\begin{center}\rule{0.5\linewidth}{0.5pt}\end{center}

\section{9. Simulation and Monte
Carlo}\label{simulation-and-monte-carlo}

\texttt{pyx.simulate\_panel(design,\ n,\ t,\ seed=…)} draws a panel from
a named design and returns it in long form (\texttt{id}, \texttt{t},
\texttt{y}, \texttt{q}, \ldots), ready for the estimators;
\texttt{pyx.monte\_carlo(design,\ n,\ t,\ …)} replicates the estimation
and summarises bias, coverage, size and power.

{\def\LTcaptype{none} % do not increment counter
\begin{longtable}[]{@{}
  >{\raggedright\arraybackslash}p{(\linewidth - 4\tabcolsep) * \real{0.3300}}
  >{\raggedright\arraybackslash}p{(\linewidth - 4\tabcolsep) * \real{0.2000}}
  >{\raggedright\arraybackslash}p{(\linewidth - 4\tabcolsep) * \real{0.4700}}@{}}
\toprule\noalign{}
\begin{minipage}[b]{\linewidth}\raggedright
Design
\end{minipage} & \begin{minipage}[b]{\linewidth}\raggedright
Source
\end{minipage} & \begin{minipage}[b]{\linewidth}\raggedright
What it is
\end{minipage} \\
\midrule\noalign{}
\endhead
\bottomrule\noalign{}
\endlastfoot
\texttt{gong\_seo} & Gong and Seo (2026), §5 & \texttt{q} is an AR(1)
correlated with the errors, predetermined \\
\texttt{gong\_seo\_endog} & Gong and Seo (2026) & the same with
\texttt{q} weakly endogenous \\
\texttt{linear} & Gong and Seo (2026) & no threshold effect: the design
for the \textbf{size} of the linearity test \\
\texttt{seo\_shin\_jump}, \texttt{seo\_shin\_cont} & Seo and Shin
(2016), §6 & self-exciting, \texttt{q\ =\ y(t-1)}; discontinuous (eq.
16) and continuous (eq. 17) \\
\texttt{seo\_shin\_jump\_common}, \texttt{seo\_shin\_cont\_common} & Seo
and Shin (2016), §6.1 & intercept fixed across regimes \\
\texttt{seo\_shin\_exogenous}, \texttt{seo\_shin\_exogenous\_common} &
Seo and Shin (2016), App. A.4 &
\texttt{q\ \textasciitilde{}\ U{[}-1,\ 1{]}}, independent of the
errors \\
\end{longtable}
}

\textbf{\texttt{simulate\_panel(design,\ n,\ t,\ *,\ seed,\ burn=50,\ unbalanced=0.0,\ min\_obs=3,\ **changes)}}
returns a frame with the columns \texttt{id}, \texttt{t} (1 to
\texttt{t}), \texttt{y} and \texttt{q}; the panel is a function of
\texttt{seed} alone. \texttt{burn} periods are simulated and discarded
first, so that \texttt{q} and \texttt{y} are close to their stationary
behaviour. \texttt{unbalanced} is the share of observations dropped at
random (the remaining values do not change), and a unit left with fewer
than \texttt{min\_obs} observations is removed. \texttt{**changes}
override any parameter of the design (\texttt{gamma}, \texttt{sigma},
\texttt{rho}, \texttt{rho\_eu}, \texttt{beta\_y}, \texttt{delta0},
\ldots) and \texttt{jump=} sets the jump of the regression function at
the threshold (\texttt{0} is a kink).
\texttt{resolve\_design(name,\ **changes)} returns the resulting
\texttt{Design}: its \texttt{truth} holds the true coefficients (named
like the estimates, plus \texttt{gamma}) and \texttt{arguments()} the
call that estimates it (the roles of the variables and
\texttt{collapse}).

\textbf{\texttt{monte\_carlo(design,\ n,\ t,\ *,\ reps=200,\ seed,\ workers,\ alpha=0.05,\ shifts=(0.10,\ 0.25,\ 0.50),\ simulation=None,\ **options)}}
replicates the estimation. \texttt{simulation} is a dictionary of
generator parameters (\texttt{jump}, \texttt{unbalanced}, \texttt{burn},
any parameter of the design); \texttt{options} are \texttt{xtdpthresh}
options (\texttt{n\_boot}, \texttt{grid}, \texttt{grid\_ci},
\texttt{level}, \ldots) and override the design\textquotesingle s
arguments; \texttt{alpha} is the level of the linearity test;
\texttt{shifts} are the \texttt{c} of the alternatives \texttt{γ\ +\ c}
whose rejection rate is reported. The result keeps one row per
replication (\texttt{to\_dataframe()}, with \texttt{status} and
\texttt{error} for the replications that failed) and answers:

{\def\LTcaptype{none} % do not increment counter
\begin{longtable}[]{@{}
  >{\raggedright\arraybackslash}p{(\linewidth - 2\tabcolsep) * \real{0.2400}}
  >{\raggedright\arraybackslash}p{(\linewidth - 2\tabcolsep) * \real{0.7600}}@{}}
\toprule\noalign{}
\begin{minipage}[b]{\linewidth}\raggedright
Method
\end{minipage} & \begin{minipage}[b]{\linewidth}\raggedright
Returns
\end{minipage} \\
\midrule\noalign{}
\endhead
\bottomrule\noalign{}
\endlastfoot
\texttt{estimates()} & bias, RMSE and coverage of the threshold and of
every slope, by parameter \\
\texttt{threshold\_coverage()} & coverage of the true threshold by the
confidence set (segments) and by its hull, and the share of incomplete
sets \\
\texttt{linearity()} & rejection rate of the linearity test: the power
with a threshold effect, the size under design \texttt{linear} \\
\texttt{power()} & rejection rate of \texttt{γ\ =\ γ₀\ +\ c},
\texttt{c\ =\ 0} included (the size of the threshold test) \\
\texttt{summary()} & a short text report \\
\end{longtable}
}

\begin{Shaded}
\begin{Highlighting}[]
\NormalTok{panel }\OperatorTok{=}\NormalTok{ pyx.simulate\_panel(}\StringTok{"gong\_seo"}\NormalTok{, }\DecValTok{200}\NormalTok{, }\DecValTok{6}\NormalTok{, seed}\OperatorTok{=}\DecValTok{3}\NormalTok{)}
\BuiltInTok{print}\NormalTok{(panel.head(}\DecValTok{3}\NormalTok{))}

\ImportTok{from}\NormalTok{ pyxtdpthresh.extras.simulation }\ImportTok{import}\NormalTok{ resolve\_design}

\NormalTok{design }\OperatorTok{=}\NormalTok{ resolve\_design(}\StringTok{"gong\_seo"}\NormalTok{, gamma}\OperatorTok{=}\FloatTok{0.5}\NormalTok{, jump}\OperatorTok{=}\FloatTok{0.2}\NormalTok{)}
\NormalTok{gaps }\OperatorTok{=}\NormalTok{ pyx.simulate\_panel(design, }\DecValTok{100}\NormalTok{, }\DecValTok{6}\NormalTok{, seed}\OperatorTok{=}\DecValTok{3}\NormalTok{, burn}\OperatorTok{=}\DecValTok{20}\NormalTok{, unbalanced}\OperatorTok{=}\FloatTok{0.15}\NormalTok{, min\_obs}\OperatorTok{=}\DecValTok{3}\NormalTok{)}
\BuiltInTok{print}\NormalTok{(}\BuiltInTok{len}\NormalTok{(gaps), gaps[}\StringTok{"id"}\NormalTok{].nunique(), design.truth[}\StringTok{"gamma"}\NormalTok{], }\BuiltInTok{round}\NormalTok{(design.jump, }\DecValTok{2}\NormalTok{))}
\NormalTok{gap\_fit }\OperatorTok{=}\NormalTok{ pyx.xtdpthresh(gaps, }\StringTok{"y"}\NormalTok{, }\BuiltInTok{id}\OperatorTok{=}\StringTok{"id"}\NormalTok{, time}\OperatorTok{=}\StringTok{"t"}\NormalTok{, bootstrap}\OperatorTok{=}\VariableTok{False}\NormalTok{, }\OperatorTok{**}\NormalTok{design.arguments())}
\BuiltInTok{print}\NormalTok{(}\BuiltInTok{round}\NormalTok{(gap\_fit.gamma, }\DecValTok{3}\NormalTok{), gap\_fit.n\_obs)}

\NormalTok{mc }\OperatorTok{=}\NormalTok{ pyx.monte\_carlo(}\StringTok{"gong\_seo"}\NormalTok{, }\DecValTok{100}\NormalTok{, }\DecValTok{5}\NormalTok{, reps}\OperatorTok{=}\DecValTok{4}\NormalTok{, seed}\OperatorTok{=}\DecValTok{11}\NormalTok{, workers}\OperatorTok{=}\DecValTok{1}\NormalTok{, n\_boot}\OperatorTok{=}\DecValTok{49}\NormalTok{, grid\_ci}\OperatorTok{=}\DecValTok{20}\NormalTok{)}
\BuiltInTok{print}\NormalTok{(mc.threshold\_coverage())     }\CommentTok{\# coverage of the true threshold by the set and by its hull}
\BuiltInTok{print}\NormalTok{(mc.linearity())              }\CommentTok{\# rejection rate: the power here, the size under design "linear"}
\BuiltInTok{print}\NormalTok{(mc.power())                  }\CommentTok{\# rejection rate of gamma + c, c = 0 included}

\NormalTok{mc2 }\OperatorTok{=}\NormalTok{ pyx.monte\_carlo(}\StringTok{"gong\_seo"}\NormalTok{, }\DecValTok{100}\NormalTok{, }\DecValTok{5}\NormalTok{, reps}\OperatorTok{=}\DecValTok{4}\NormalTok{, seed}\OperatorTok{=}\DecValTok{11}\NormalTok{, workers}\OperatorTok{=}\DecValTok{1}\NormalTok{, n\_boot}\OperatorTok{=}\DecValTok{49}\NormalTok{, grid\_ci}\OperatorTok{=}\DecValTok{20}\NormalTok{,}
\NormalTok{                      alpha}\OperatorTok{=}\FloatTok{0.10}\NormalTok{, shifts}\OperatorTok{=}\NormalTok{(}\FloatTok{0.2}\NormalTok{, }\FloatTok{0.4}\NormalTok{), simulation}\OperatorTok{=}\NormalTok{\{}\StringTok{"jump"}\NormalTok{: }\FloatTok{0.5}\NormalTok{, }\StringTok{"unbalanced"}\NormalTok{: }\FloatTok{0.1}\NormalTok{\})}
\BuiltInTok{print}\NormalTok{(mc2.power())}
\BuiltInTok{print}\NormalTok{(mc2.estimates().}\BuiltInTok{round}\NormalTok{(}\DecValTok{3}\NormalTok{))}
\end{Highlighting}
\end{Shaded}

The coverage and power figures of four replications mean nothing; the
tables in \texttt{docs/montecarlo/} use hundreds. The calls are the
same.

The experiment is a function of \texttt{seed} alone, whatever
\texttt{workers} is. \texttt{workers=None} uses
\texttt{min(cpu\_count,\ reps)} processes of one thread each;
\texttt{workers=1} runs in the calling process.

\textbf{On Windows, where worker processes are started by spawning, a
script that runs replications in several processes must be a real file
with the usual guard:}

\begin{Shaded}
\begin{Highlighting}[]
\CommentTok{\# not run by the tests}
\ImportTok{import}\NormalTok{ pyxtdpthresh }\ImportTok{as}\NormalTok{ pyx}

\ControlFlowTok{if} \VariableTok{\_\_name\_\_} \OperatorTok{==} \StringTok{"\_\_main\_\_"}\NormalTok{:}
\NormalTok{    mc }\OperatorTok{=}\NormalTok{ pyx.monte\_carlo(}\StringTok{"gong\_seo"}\NormalTok{, }\DecValTok{400}\NormalTok{, }\DecValTok{6}\NormalTok{, reps}\OperatorTok{=}\DecValTok{200}\NormalTok{, seed}\OperatorTok{=}\DecValTok{1}\NormalTok{, workers}\OperatorTok{=}\DecValTok{4}\NormalTok{)}
    \BuiltInTok{print}\NormalTok{(mc.summary())}
\end{Highlighting}
\end{Shaded}

The worker processes re-import the main file; without the guard each one
would start its own experiment, and a script read from standard input
(\texttt{python\ -}) has no file to re-import: the replications fail.
\texttt{examples/04\_monte\_carlo.py} is a complete case.

The bundled results of larger experiments, and what they say about the
agreement with the papers (coverage agrees; the power to reject shifted
thresholds is lower than Table 2 of Gong and Seo, cause not
established), are in \texttt{docs/montecarlo/}.

\begin{center}\rule{0.5\linewidth}{0.5pt}\end{center}

\section{10. Errors and warnings}\label{errors-and-warnings}

\textbf{Exceptions} all derive from \texttt{PyxdpthreshError}:

{\def\LTcaptype{none} % do not increment counter
\begin{longtable}[]{@{}
  >{\raggedright\arraybackslash}p{(\linewidth - 2\tabcolsep) * \real{0.3800}}
  >{\raggedright\arraybackslash}p{(\linewidth - 2\tabcolsep) * \real{0.6200}}@{}}
\toprule\noalign{}
\begin{minipage}[b]{\linewidth}\raggedright
Class
\end{minipage} & \begin{minipage}[b]{\linewidth}\raggedright
Raised when
\end{minipage} \\
\midrule\noalign{}
\endhead
\bottomrule\noalign{}
\endlastfoot
\texttt{SpecificationError} (also a \texttt{ValueError}) & the call
breaks a rule of the Stata command; \texttt{.stata\_code} is
Stata\textquotesingle s return code (198 for most option errors) and
\texttt{.key} names the rule \\
\texttt{DataError} (also a \texttt{ValueError}) & the data cannot be
used: repeated (unit, time) pairs, missing keys, no usable observation,
an unbalanced panel for \texttt{xthenreg}, a point estimation that finds
no valid solution \\
\texttt{NumericalError} & a computation fails numerically \\
\end{longtable}
}

All of them derive from \texttt{pyx.PyxdpthreshError} and are importable
from the package (\texttt{pyx.SpecificationError},
\texttt{pyx.DataError}, \texttt{pyx.NumericalError}).
\texttt{SpecificationError} and \texttt{DataError} are also
\texttt{ValueError}s.

\begin{Shaded}
\begin{Highlighting}[]
\ControlFlowTok{try}\NormalTok{:}
\NormalTok{    pyx.xtdpthresh(df, }\StringTok{"invest"}\NormalTok{, grid}\OperatorTok{=}\DecValTok{5}\NormalTok{, bootstrap}\OperatorTok{=}\VariableTok{False}\NormalTok{, }\OperatorTok{**}\NormalTok{MODEL)}
\ControlFlowTok{except}\NormalTok{ pyx.SpecificationError }\ImportTok{as}\NormalTok{ error:}
    \BuiltInTok{print}\NormalTok{(error.stata\_code, error.key)         }\CommentTok{\# 198 opt\_grid}
    \BuiltInTok{print}\NormalTok{(error)}
\end{Highlighting}
\end{Shaded}

The return codes match the command\textquotesingle s. On 19 hostile
inputs (minimal panels, a constant, binary or rare threshold variable,
an empty regime, duplicated regressors, fewer units than instruments, a
null outcome) the package returns the same code as Stata in all 19.

\textbf{"No threshold admitted a valid GMM solve"} (a
\texttt{DataError}) means that at every candidate threshold the
estimator could not be computed. The usual causes are too few units or
periods, more instruments than the data can support, a threshold
variable with too little variation, or a regime that is too small.
Remedies: \texttt{collapse=True}, a tighter \texttt{maxlag}, a larger
\texttt{trim}.

\textbf{Warnings} are \texttt{pyx.PyxdpthreshWarning} (a
\texttt{UserWarning}), so the \texttt{warnings} module controls them,
and \texttt{warn=False} silences most of them at the source. They cover:
many instrument columns, memory, a kink with \texttt{q} not among the
regressors, \texttt{refine} not applied after a one-step fallback, a
Windmeijer correction that could not be applied, and the legacy
\texttt{coef\_boot}.

\begin{center}\rule{0.5\linewidth}{0.5pt}\end{center}

\section{11. What differs from Stata}\label{what-differs-from-stata}

{\def\LTcaptype{none} % do not increment counter
\begin{longtable}[]{@{}
  >{\raggedright\arraybackslash}p{(\linewidth - 2\tabcolsep) * \real{0.4200}}
  >{\raggedright\arraybackslash}p{(\linewidth - 2\tabcolsep) * \real{0.5800}}@{}}
\toprule\noalign{}
\begin{minipage}[b]{\linewidth}\raggedright
Stata
\end{minipage} & \begin{minipage}[b]{\linewidth}\raggedright
Python
\end{minipage} \\
\midrule\noalign{}
\endhead
\bottomrule\noalign{}
\endlastfoot
\texttt{xtset\ firm\ year} & \texttt{id="firm",\ time="year"}, or a
\texttt{MultiIndex}, or \texttt{from\_stata} \\
\texttt{if} / \texttt{in} & \texttt{where=}; no \texttt{in} \\
\texttt{e(b)}, \texttt{e(V)}, \texttt{e(N)}, \ldots{} &
\texttt{res.params}, \texttt{res.cov}, \texttt{res.scalars{[}"N"{]}},
\ldots{} \\
\texttt{ereturn\ list,\ all} (hidden results) & \texttt{res.scalars}
includes the hidden ones that are kept for checking \\
\texttt{predict,\ residuals\ \textbackslash{}\textbar{}\ arresiduals\ \textbackslash{}\textbar{}\ xb\ \textbackslash{}\textbar{}\ regime}
&
\texttt{res.predict("residuals"\ \textbackslash{}\textbar{}\ "ar\_residuals"\ \textbackslash{}\textbar{}\ "xb"\ \textbackslash{}\textbar{}\ "regime",\ frame)} \\
\texttt{estimates\ store} / \texttt{esttab} & the result objects,
\texttt{pyx.compare({[}...{]})} and the \texttt{to\_*} methods \\
\texttt{set\ level} & \texttt{level=} of the call (95 by default) \\
\texttt{rseed(\#)}, not restored after the command & \texttt{seed=}; no
global state is changed \\
\texttt{verbose}, \texttt{exportgmm} & accepted with no effect \\
\texttt{xthenreg\ …\ h\_0()} & \texttt{xthenreg(…,\ h\_0=)} \\
a missing value in \texttt{if} counts as true & counts as false \\
Mata matrices left in memory & frozen containers under
\texttt{res.fit} \\
\end{longtable}
}

What the package adds: the result classes with tables, plots and a JSON
record; \texttt{diagnose()} and \texttt{settings\_log()};
\texttt{robustness()}; the simulation toolkit; bundled data; English and
French messages. What it does not change: the estimators, the bootstrap
draws, the tests and their conventions, which are checked against Stata
17 MP to four decimals.

% Methodological note of pyxtdpthresh: Appendix A of the user guide.
% Inserted before the end of the document by tools/build_guide_pdf.py, which also replaces the
% line that includes ported_functions.tex by the tables generated from docs/ported_functions.md.
\clearpage
\setcounter{secnumdepth}{3}
\appendix
\providecommand{\E}{\operatorname{E}}
\newcommand{\code}[1]{\texttt{\detokenize{#1}}}
\newcommand{\ind}{\mathbf{1}}
\newcommand{\bZ}{\mathbf{Z}}
\newcommand{\bX}{\mathbf{X}}
\newcommand{\diff}{\Delta}
\newcommand{\statasrc}[1]{\texttt{\detokenize{#1}}}

\section{Methodological note}\label{app:note}

This appendix is the methodological note of the package. It states what the two estimators
compute (Sections~\ref{note:model}--\ref{note:var}), how the package computes it and why it is
fast (Section~\ref{note:cumsum}), how the bootstrap inference is organised
(Section~\ref{note:boot}), where it departs from the algorithm of Gong and Seo
(Section~\ref{note:gs}), what \texttt{xthenreg} adds (Section~\ref{note:xthenreg}), which conventions of the
Stata commands are reproduced on purpose (Section~\ref{note:quirks}), what has been compared with Stata and how
(Section~\ref{note:valid}), and where every Stata function went (Section~\ref{note:map}).

Three kinds of statements are mixed and kept apart in the text. A \emph{model} statement comes from
the papers the commands implement (Seo and Shin 2016; Seo, Kim and Kim 2019; Gong and Seo 2026;
Nguyen and Lai 2026). A \emph{convention} is a behaviour of the Stata code that the package
reproduces so that the numbers agree; Section~\ref{note:quirks} lists them. An \emph{implementation}
statement says how the package obtains a result and is checked by tests, not by a theorem.
Nothing here claims more theory than the papers prove: where a scheme is an approximation that the
theory does not cover, the text says so.

\subsection{Notation}\label{note:notation}

\begin{longtable}{@{}p{2.6cm}p{13.4cm}@{}}
\toprule
Symbol & Meaning \\
\midrule
\endhead
$i=1,\dots,n$ & units (the clusters of the variance and of the bootstrap); $n$ is \code{e(N_units)} \\
$e$ & one transformed equation: a unit and a period. There are $N$ of them, \code{e(N)}, stacked by unit then period \\
$y_{it}, x_{it}, q_{it}$ & dependent variable, the $k$ regressors (the lag of $y$ first, in the dynamic model) and the threshold variable \\
$\gamma$ & the threshold; $\Gamma$ is the grid of candidates, $G$ its size \\
$\mathcal T$ & the transformation (first differences or forward orthogonal deviations) as a sparse matrix that maps the \emph{level} observations to the $N$ equations \\
$\theta=(\beta',\delta')'$ & coefficients; $p=2k+1$ in the jump model, $p=k+1$ in the kink model \\
$\bX(\gamma)$ & transformed regressors at $\gamma$: $\mathcal T\,[\,x,\ \ind(q>\gamma),\ x\,\ind(q>\gamma)\,]$ (jump) or $\mathcal T\,[\,x,\ (q-\gamma)\ind(q>\gamma)\,]$ (kink) \\
$\bZ$ & the $N\times L$ instrument matrix; $z_e$ is its row $e$ \\
$D(\gamma)$, $m_y$ & $D(\gamma)=\bZ'\bX(\gamma)/N$ ($L\times p$) and $m_y=\bZ'\,\mathcal T y/N$ ($L$) \\
$g_i(\theta,\gamma)$ & moment contribution of unit $i$: $\sum_{e\in i} z_e\,(\mathcal T y_e-\bX_e(\gamma)'\theta)$ \\
$W_1,\ W_2$ & weight matrices of the first and second step \\
$Q(\theta,\gamma)$ & criterion $N\,\bar g'\,W\,\bar g$ with $\bar g=m_y-D(\gamma)\theta$ \\
\bottomrule
\end{longtable}

Stata counts the stacked equations $N$ in \code{e(N)} and the units in \code{e(N_units)}; the note
follows that. (The papers write $n$ for the units and index the moments by period.)

\subsection{Model and moment conditions}\label{note:model}

The commands fit the single-threshold dynamic panel model of Seo and Shin (2016),
\begin{equation}\label{eq:jump}
y_{it}=x_{it}'\beta+(1,x_{it}')\,\delta\;\ind(q_{it}>\gamma)+\eta_i+\varepsilon_{it},
\end{equation}
where $x_{it}$ contains $y_{i,t-1}$ unless the model is static. The regression function jumps at
$\gamma$ by $\delta_1+x'\delta_2$. With \code{kink}, continuity is imposed,
\begin{equation}\label{eq:kink}
y_{it}=x_{it}'\beta+\delta_q\,(q_{it}-\gamma)\,\ind(q_{it}>\gamma)+\eta_i+\varepsilon_{it},
\end{equation}
and $q$ should be a regressor so that it has a slope below $\gamma$. The unit effect $\eta_i$ is
removed by a transformation.

\paragraph{First differences.} The equation of unit $i$ at period $t$ exists when the unit has rows
at the calendar periods $t$ and $t-1$ (a gap in the dates gives no equation) and both are in the
equation sample and complete in $y$, $q$ and the regressors. It is $y_t-y_{t-1}$ on
$x_t-x_{t-1}$, and touches two level observations with weights $+1$ and $-1$.

\paragraph{Forward orthogonal deviations.} An equation row followed by $T_f\ge1$ later equation rows of
the same unit is $c\,(y_t-\overline{y}_{\text{later}})$ with $c=\sqrt{T_f/(T_f+1)}$ and the same for
every regressor; the dates of the later rows do not matter, and the last row of a unit has no
equation. In weights: $+c$ on the row, $-c/T_f$ on each later row.

\paragraph{Triplets.} Both transformations are stored as triplets (equation, level observation,
weight), i.e.\ as a sparse matrix $\mathcal T$ with two entries per row for first differences. The
regime regressors depend on $\gamma$ through the level observations only, so $\bX(\gamma)$ is the
transform of the level regressors $[x,\ \ind(q>\gamma),\ x\,\ind(q>\gamma)]$: the regressors at any
threshold come from one sparse product, and the equations are never restacked. This is what makes
the cumulative-sum algorithm of Section~\ref{note:cumsum} possible.

\paragraph{Instruments.} The vector $z_e$ of an equation at period $t$ contains (Arellano and Bond 1991;
Roodman 2009), each in its own block of columns per period unless \code{collapse}:
a constant; under first differences the levels $y_{t-2},y_{t-3},\dots$; the transformed strictly
exogenous regressors; the lagged levels of the endogenous regressors from $t-2$ and of the
predetermined regressors from $t-1$; and the external instruments dated $t$. Under forward
orthogonal deviations the dates are one period later, since the transformed error involves later
periods. \code{maxlag(lo hi)} bounds the lags, the minimum being 2 for $y$ and the endogenous
variables and 1 for the predetermined ones. With \code{collapse} the per-period blocks are replaced
by one shared block (one column per lag instead of one per lag and period). An equation without an
instrument other than the constant is dropped. The instrument matrix does not depend on $\gamma$.

Columns that carry no information are then removed, in the order of the Stata command:
(1)~columns that are multiples of the period constants (a variable common to all units in a period)
are zeroed; (2)~columns that are exactly zero are dropped; (3)~exact duplicates are dropped, the
first kept; (4)~columns that are linear combinations of the earlier ones are dropped (squared
relative residual below $10^{-13}$). The counts are \code{e(N_iv_common)}, \code{e(N_iv_dep)} and
\code{e(N_iv_dep_near)}.

\paragraph{Moment conditions.} With $T(\cdot)$ the chosen transformation, the estimator uses
$\E[z_e\,\mathcal T(\varepsilon)_e]=0$. The sample moment is
$\bar g(\theta,\gamma)=m_y-D(\gamma)\theta$ and the criterion is
$Q(\theta,\gamma)=N\,\bar g'W\bar g$.

\paragraph{Time effects.} With \code{td} the transformed period dummies are partialled out of the
transformed $y$, regressors and instruments (Frisch--Waugh--Lovell). Under first differences the
transformed dummy of a period is the same for every unit, so the projection is exact demeaning within
the period; a period with one equation is demeaned to zero and its row is dropped. Under forward
orthogonal deviations the dummies receive the same transform as the data and the system is projected
on them with a generalised inverse (\code{invsym}); a row that the projection saturates is dropped.
The projections $M$ are symmetric and idempotent, so $(M\bZ)'\,d\bX=\bZ'\,M\,d\bX$ and the regime
moments are formed from the projected instruments and the unprojected level regressors.

\subsection{Estimation}\label{note:est}

\paragraph{Closed form at a threshold.} For a weight matrix $W$ and a given $\gamma$ the criterion is
quadratic in $\theta$, so
\begin{equation}\label{eq:theta}
\hat\theta(\gamma)=\bigl(D'WD\bigr)^{-1}D'W\,m_y,\qquad
\bar g=m_y-D\hat\theta,\qquad Q=N\,\bar g'W\bar g .
\end{equation}
The inverse is Mata's \code{invsym}, or the gated \code{syminv} of \code{xtdpthresh}
(Section~\ref{note:quirks}): a singular normal matrix is not an error but a candidate that is not
admissible.

\paragraph{First step.} $\theta_1$ and $\gamma_1$ minimise $Q$ over the grid with $W_1$. Under first
differences $W_1$ is the MA(1) weight of Arellano and Bond (1991): the transformed errors of
consecutive equations of a unit overlap in one term, so with $S_0=\bZ'\bZ/n$ and
$S_1=n^{-1}\sum_i\sum_t z_{i,t-1}z_{it}'$ over consecutive equations of the same unit,
\begin{equation}\label{eq:w1}
W_1=\bigl(2S_0-S_1-S_1'\bigr)^{-1}.
\end{equation}
If that matrix cannot be inverted, $(\bZ'\bZ)^{-1}$ is used, then the identity
(\code{e(W1_fallback)} $=1$ or $2$). Under forward orthogonal deviations $W_1=(\bZ'\bZ/N)^{-1}$.

\paragraph{Second step.} With the first-step residuals $r_1$ at $(\theta_1,\gamma_1)$, the cluster
covariance of the moments is
\begin{equation}\label{eq:omega}
\Omega=\frac1N\sum_{i=1}^n g_ig_i'-\frac{1}{nN}\Bigl(\sum_i g_i\Bigr)\Bigl(\sum_i g_i\Bigr)',
\end{equation}
the second term being dropped under \code{nocenter}. $W_2=\Omega^{-1}$ is built \emph{once}, from the
first-step estimate, and then \emph{held fixed} while the second-step criterion is evaluated at every
other threshold; $\hat\gamma$ minimises that criterion over the grid, with ties going to the smaller
threshold (relative tolerance $10^{-12}$). This is what the Stata search does, and it is why the
criterion of every candidate can be computed from sufficient statistics.

\paragraph{One-step fallback.} If $\Omega$ cannot be inverted, no candidate can be solved with $W_2$,
or the criterion does not vary over the admitted candidates of the second step, the first-step
estimates are reported with the sandwich of $W_1$ (\code{e(estimator_twostep)} $=0$) and without
Hansen's $J$. If the first-step criterion itself is flat, the command stops with return code 498.

\paragraph{The grid.} The effective support of $q$ is the set of $q$ values of the level observations
that enter a stacked equation (periods $t$ and $t-1$ under first differences); \code{gridsample(observed)}
keeps the current-row values only. The trimming bounds are the type-7 quantiles at
\code{trim}/2 and $1-$\code{trim}/2. A uniform grid has equally spaced points between the bounds;
a quantile grid places them on the empirical quantiles (duplicates collapsed). A candidate is
admitted when both regimes have at least $m$ observations, where the default floor is
$m=\max\{2,\min(\#\{q\le q_{lo}\},\#\{q>q_{hi}\})\}$ and \code{minregime} can only raise it, and when
there are more instruments than parameters. At least 20 equations, 5 units and 20 usable values of
$q$ are required. The grids of 100, 199 and 397 uniform points are nested ($g\to2g-1\to4g-3$), so
a finer grid is a direct check.

\paragraph{Refinement.} Under \code{kink} the criterion is continuous in $\gamma$, and $\hat\gamma$
minimises it between the two grid points next to the grid minimum: up to 8 rounds of 40 and then 10
trial points, shrinking the bracket until its spacing is below $10^{-6}$ of its width
(\code{e(kink_refined)}). Under the jump model \texttt{refine(\#)} adds the observed values of $q$
between the neighbours of the minimiser, in at most \# rounds, and re-minimises the second-step
criterion with the same fixed $W_2$; the confidence set and the tests keep using the initial grid.

\subsection{Cumulative sums: the criterion on a whole grid in one pass}\label{note:cumsum}

Only the regime columns of $\bX(\gamma)$ depend on $\gamma$. For instrument $j$ and level regressor
$k$ the regime moment is
\begin{equation}\label{eq:S}
S_{jk}(\gamma)=\sum_{e}\bZ_{ej}\sum_{(l,w)\in e} w\,W_{lk}\,\ind(q_l>\gamma),
\end{equation}
where the inner sum runs over the level observations $l$ of equation $e$ with their weights $w$
(for the jump model $W_l=(1,x_l)$). Sort the grid, $\gamma_1<\dots<\gamma_G$, and let
$m_l=\#\{g:\gamma_g<q_l\}$. Then $\ind(q_l>\gamma_g)=1$ exactly when $g\le m_l$, so each term
$c=\bZ_{ej}\,w\,W_{lk}$ contributes to the first $m_l$ grid points and to no other. The kernel adds
each contribution \emph{once}, to the bucket $m_l$, and a reverse cumulative sum over the buckets
gives every grid point:
\begin{equation}\label{eq:cumsum}
A[m]=\sum_{\{p:\,m_p=m\}}c_p,\qquad S(\gamma_g)=\sum_{m\ge g}A[m].
\end{equation}
The cost is $O(P\cdot\mathrm{nnz}\cdot K+G\,L\,K)$, with $P$ the number of pairs and $\mathrm{nnz}$
the number of non-zero entries of an instrument column, against $O(G\,P\,\mathrm{nnz}\,K)$ if each
threshold were recomputed. The kernel runs over instruments (one Numba thread owns the slice
$S[:,j,:]$), reads $\bZ$ in compressed-column form so that the block-sparse layout costs nothing, and
places a missing $q$ in bucket 0, so that it is never above the threshold.

For the kink model the regime column is $(q-\gamma)\ind(q>\gamma)=q\,\ind(q>\gamma)-\gamma\,\ind(q>\gamma)$,
so the same pass with level columns $(1,q)$ gives $S_\kappa(\gamma)=S_q(\gamma)-\gamma\,S_1(\gamma)$.
The full matrix is $D(\gamma_g)=[\,\bZ'\bX_{\text{lin}}/N\ \big|\ S(\gamma_g)/N\,]$, the first block
being the same at every threshold.

\paragraph{Solving the whole grid.} Given $D(\gamma_g)$ for all $g$, the estimates (\ref{eq:theta}) are
batched matrix products (BLAS) and one batched \code{invsym}; this step is dominated by dense algebra
and is plain NumPy. On the $565\times15$ panel the regime moments of 100 grid points take 4~ms
against 278~ms for the straightforward reference implementation that is kept for the equivalence
tests (block~C of the project plan), and the two agree to $10^{-12}$. At realistic sizes the
remaining time is in the bootstrap (Section~\ref{note:boot}), whose criterion is one matrix product
that already runs close to the speed of the machine, which is why no further algebraic gain exists
there: an idea of sharing the unrestricted fit across candidates was examined and rejected because the
residual depends on the candidate.

\paragraph{Unit moments.} $g_i=\sum_{e\in i}z_e r_e$ is a segmented sum over the CSR layout of $\bZ$
(one kernel). Units without equations are zero rows and change nothing in (\ref{eq:omega}).

\subsection{Variances and specification tests}\label{note:var}

\paragraph{Robust sandwich.} For a fixed weight $A$ (efficient or not), with $D=\bZ'\bX/N$ and
$B=(D'AD)^{-1}$,
\begin{equation}\label{eq:sandwich}
\widehat V=\frac1N\,B\,D'A\,\Omega\,A D\,B,
\end{equation}
with $\Omega$ rebuilt from the residuals at the reported estimate. It is valid for the one-step
estimator too. This is the variance \emph{conditional on $\hat\gamma$}.

\paragraph{Joint variance with the threshold.} Following Seo and Shin (2016) the threshold is one more
parameter. The Jacobian is augmented to $[\,D\ |\ \bZ'x_\gamma/N\,]$, where $x_\gamma$ is the
derivative of the transformed residual with respect to $\gamma$, and the $\theta$ block of the
sandwich of the augmented system is reported. In the jump model the sample moment is a step function
of $\gamma$, so the derivative of its expectation is smoothed by a Gaussian kernel,
\begin{equation}\label{eq:xg}
x_\gamma=-\,\mathcal T\Bigl[\varphi_h(q-\hat\gamma)\,(\hat\delta_1+x'\hat\delta_2)\Bigr],\qquad
h=c\cdot1.06\cdot s_q\cdot n^{-1/5},
\end{equation}
with $\varphi_h(u)=\varphi(u/h)/h$, $s_q$ the standard deviation of $q$ over all level observations,
$n$ the number of units with an equation and $c=$ \code{bwscale} (1.5 by default, as in
\code{xthenreg}). In the kink model the derivative is analytic, $x_\gamma=-\hat\delta_q\,\mathcal T\,\ind(q>\hat\gamma)$.
The variance conditional on $\hat\gamma$ is kept (\code{e(V_cond)}); if the joint variance cannot be
computed, the conditional one is reported (\code{e(joint_vce)} $=0$). The upper-regime slopes are
$\hat\beta+\hat\delta_2$, with the variance of the linear combination.

\paragraph{Windmeijer correction.} $W_2=\Omega(\theta_1)^{-1}$ depends on the first-step estimate.
Linearising the second-step estimate in $\theta_1$ (Windmeijer 2005) adds
$D_j=-B_2G_2'W_2\,[\partial\Omega/\partial\theta_j]\,W_2\bar g_2$, with $B_2=(G_2'W_2G_2)^{-1}$,
$G_2=D(\hat\gamma)$, $\bar g_2$ the second-step moment residual and $B_1,G_1$ the first-step
counterparts. The variance is
\begin{equation}\label{eq:wind}
\widehat V_W=\frac1N\,M\,\Omega_1\,M',\qquad M=B_2G_2'W_2+[D_1,\dots,D_p]\;B_1G_1'W_1 .
\end{equation}
The derivative of $\Omega$ is never formed. With $h_{ij}=\sum_{e\in i}z_ex_{ej}$ and
$s=\sum_ig_i$, only the product
$[\partial\Omega/\partial\theta_j]\,v=-\frac1N\sum_i\bigl(h_{ij}\,g_i'v+g_i\,h_{ij}'v\bigr)
+\frac{1}{nN}\bigl(h_j\,s'v+s\,h_j'v\bigr)$ with $v=W_2\bar g_2$ is needed (the last term comes from the
centring), which costs two products per column instead of an $L\times L$ matrix. The correction is
computed once at the final estimate; if it cannot be computed or is not positive semidefinite, the
robust variance is reported (\code{e(vce_applied)} $=0$). Windmeijer derived it for moments linear
in the parameters; it need not improve the variance for the nonlinear moments of the threshold.

\paragraph{Hansen's $J$.} $J=Q(\hat\theta,\hat\gamma)$ of the second step, referred to $\chi^2$ with
$L-p-1$ degrees of freedom (the threshold counts as one parameter). Nothing is reported after the
one-step fallback or when the model is not over-identified.

\paragraph{Arellano--Bond serial correlation.} On the \emph{first-difference} residuals, whatever the
estimation transformation (the \code{xtabond2} convention), the statistic of order $m$ is
$\mathrm{AR}(m)=b_0/\sqrt{T_1+T_2+T_3}$ (Arellano and Bond 1991, eq.~8), where, with
$r_{it}$ the residuals and $u_i=\sum_t r_{it}r_{i,t-m}$,
$b_0=\sum_iu_i$, $T_1=\sum_iu_i^2$, $T_2=-2\,g'(MAM')^{-1}MA\,s$ and $T_3=g'\widehat V g$. Here
$g=\sum_e x^*_e\,r_{e,t-m}$ collects the regressors against the lagged residual, $s=\bZ'(r\circ u_{\text{unit}})$
the instrument moments weighted by the unit sums, $M=\bZ'\bX$ (raw sums) and $\widehat V$ the variance
of the parameters. With the joint variance the threshold column $x_\gamma$ is a parameter of $g$, $M$ and
$\widehat V$ (\code{e(ar_joint)} $=1$); the statistics that treat $\hat\gamma$ as known are
\code{e(ar1_cond)} and \code{e(ar2_cond)}. The p-value is $2\Phi(-|\mathrm{AR}|)$. At least five
residual pairs, from at least five units, are needed.

\subsection{Bootstrap inference}\label{note:boot}

\paragraph{The wild bootstrap of the command.} Every bootstrap component compares, on each sample, the
criterion of a \emph{restricted} model with the smallest criterion of the threshold models over the
initial grid. A sample is $Y^*_b=F+r\,\eta_{i(e),b}$: a fitted part $F$ and the residuals $r$ of the
restricted fit, multiplied by one weight $\eta_{ib}$ per unit (so the correlation within a unit and
the heteroskedasticity across units are kept), the same weights for every candidate. Regressors and
instruments keep their sample values. The weights follow Mammen (1993): $-\phi$ with probability
$(\sqrt5+1)/(2\sqrt5)$ and $1/\phi$ otherwise, $\phi=(\sqrt5-1)/2$; mean 0, variance 1.

\paragraph{Why it is cheap.} With a fixed weight $W$ the criterion of a candidate $g$ depends on the
sample only through the moments $m^*=\bZ'Y^*/N$:
\begin{equation}\label{eq:quad}
\theta=C_g\,m^*,\quad C_g=(D_g'WD_g)^{-1}D_g'W,\qquad
J_g(m^*)=N\,{m^*}'\,\bigl[M_g'WM_g\bigr]\,m^*,\quad M_g=I-D_gC_g .
\end{equation}
The moments of all samples are $m^*_b=\bZ'F/N+N^{-1}\sum_i\eta_{ib}\,g_i$, where $g_i=\sum_{e\in i}z_er_e$
is the unit moment of the residuals. No $N\times B$ matrix is ever formed: for $G$ candidates and $B$
draws, all the criteria are a few batched products of the $G$ matrices
$Q_g=M_g'WM_g$ ($L\times L$) with the $L\times B$ matrix of moments, in chunks that bound the memory
(8 million elements for the work array; the $Q_g$ are kept when they fit in 40 million elements and
the estimates are solved draw by draw otherwise). The bootstrap is 80--92~\% of the time of a full
run; it is evaluated this way for all draws and candidates at once.

\paragraph{Confidence set (grid bootstrap).} Following Hansen (1999b) and Gong and Seo (2026, \S4.1), the
set collects the candidates $\gamma_\ell$ that a bootstrap test does not reject. With the weight of
the reported estimator ($W_2$ after a two-step fit, $W_1$ after the fallback), for each candidate:
(1)~restricted fit $\hat\theta(\gamma_\ell)$ and residuals $e$; (2)~statistic
$D_n(\gamma_\ell)=Q(\hat\theta(\gamma_\ell),\gamma_\ell)-m_n$, where $m_n$ is the minimum of
$Q(\hat\theta(\gamma),\gamma)$ over $\Gamma\cup\{\gamma_\ell\}$ (so $D_n(\hat\gamma)=0$);
(3)~draws $Y^*_b=\bX(\gamma_\ell)\hat\theta(\gamma_\ell)+\eta_{ib}e$; (4)~$D^*_{n,b}(\gamma_\ell)=J_\ell(m^*_b)-\min\{J_\ell(m^*_b),\min_{g\in\Gamma}J_g(m^*_b)\}$;
(5)~accept $\gamma_\ell$ iff $D_n(\gamma_\ell)$ does not exceed the order statistic of rank
$k=\lceil p\,(B+1)\rceil$ of the $D^*_{n,b}$, with $p=$ \code{level}/100. That is the add-one rule:
the bootstrap p-value $(1+R)/(1+B)$, $R=\#\{D^*\ge D\}$, exceeds $1-p$. If $k>B$ no draw can reject and
every candidate is accepted. The set is the union of the accepted candidates; \code{e(ci_nseg)}
counts its segments and \code{e(gamma_lo)}, \code{e(gamma_hi)} are the hull, which contains rejected
values when there are several segments. A candidate whose restricted fit cannot be solved is not
admissible (status 3); a draw that gives no valid statistic leaves the candidate unresolved
(statuses 4--6), and then no set is reported (\code{e(ci_incomplete)} $=1$). A statistic exactly zero
is accepted mechanically (status 2).

\paragraph{Linearity test.} $H_0:\delta=0$. The statistic is the one-step criterion of the linear model
minus the smallest one-step criterion of the threshold model over the initial grid, both with $W_1$;
samples are drawn from the linear fit; the p-value is $(1+R)/(1+B)$, and is missing if any draw is invalid.
This GMM distance test differs from the sup-Wald test of \code{xthenreg}.

\paragraph{Test of a given threshold.} \texttt{citest(\#)} is one row of the grid bootstrap at $\gamma=\#$,
with its own seed, so asking for it changes no other result; its p-value is $(1+R)/(1+B)$ and its
decision is $D\le$ crit.

\paragraph{Continuity test.} Following Gong and Seo (2026, \S3.2 and \S4.3) with
$\tilde\theta=\arg\min_{\Theta_c}\hat Q_n$ the continuity-restricted estimator, the statistic is the
distance $T_n=N(Q(\tilde\theta)-Q(\hat\theta))$ between the kink and the jump criteria. The package computes
it over the initial grid with the common weight, draws samples from the kink fit with the \emph{jump}
residuals reweighted by unit, and takes the p-value $(1+R)/(1+B)$ of $T^*_n=\min_{\text{kink}}J^*-\min_{\text{jump}}J^*$.
It needs $q$ as a contemporaneous regressor, otherwise the kink model is not nested in the jump model.

\paragraph{Coefficient bootstrap (legacy).} \code{coef_boot} re-runs the whole search in each draw: samples
$Y^*=\bX(\hat\gamma)\hat\theta+r\eta$, the threshold searched again over the estimation grid with $W_1$, $\theta$
re-estimated there (two-step: $W_2^*$ from the draw's residuals, threshold searched again). The intervals are
$\hat\theta_j\pm q_{1-\alpha}(|\theta^*_j-\hat\theta_j|)$ (symmetric) or the percentiles
$q_{\alpha/2},q_{1-\alpha/2}$ of $\theta^*_j$. Section~\ref{note:gs} says why this is not the coefficient
bootstrap of Gong and Seo.

\paragraph{Unit resampling.} \code{boot_type="unit"} follows the structure of Algorithm 1
(Section~\ref{note:gs}). A resample is a vector of unit counts, so every quantity is a sum over units of
per-unit cross products that do not depend on the draw: $\bZ'\bX(\gamma_g)$ of each grid point, $\bZ'y$ and the
unit moments; a draw costs a few small products and no pass over the $N$ equations. Per draw: the
first-step coefficients of every grid point with $W_1$ of the sample, the best one giving the residuals whose
recentred cluster covariance defines $W_2^*=\Omega^{*-1}$; then the criterion at the candidate (restricted)
and at every grid point (unrestricted) with $W_2^*$. A draw whose profile has no usable minimum is skipped,
not replaced.

\paragraph{Random numbers.} Stata 17 uses the 64-bit Mersenne Twister; \code{set seed s} starts it with
\code{init_genrand64(s)}, and a draw is $u=((x\gg11)\,|\,1)\cdot2^{-53}$, where $x$ is the next 64-bit output
(measured on 10\,000 draws from five seeds: identical). A matrix is filled row by row;
\code{uniform()}, \code{runiform()} and \code{invnormal(uniform())} consume one stream. The package reproduces
the stream, so the same seed gives the same weights as the command. With \code{rseed(s)} each component restarts
at $(s+\text{offset})\bmod2^{31}$: confidence set $0$, linearity $104\,729$, continuity $224\,737$, coefficient
$350\,377$, \code{citest} $477\,377$. The weights of a component are drawn replication after replication,
one uniform per contributing unit. Without a seed one random seed is drawn and the components share a
stream in the order they run. The Python wrapper always records a seed (Section~\ref{note:quirks}).

\subsection{Gong and Seo's Algorithm 1 and what the package does}\label{note:gs}

This section is the reading of Gong and Seo (2026), Algorithm~1 and \S\S4.1--4.3 (Theorems 5--7), against the code.
The proofs were not re-derived; the comparison is of the algorithms, and the statements about validity are the
paper's.

\paragraph{The algorithm of the paper.} For a parameter $\theta^*_0=(\beta^{*\prime}_0,\delta^{*\prime}_0,\gamma^*_0)'$ that
generates the bootstrap data:
\begin{enumerate}
\item draw $i^*$ i.i.d.\ uniformly from $\{1,\dots,n\}$ and resample the units \emph{jointly}:
$(x^*_{it},x^*_{it-1},z^*_{it},\widehat{\Delta\epsilon}^*_{it})=(x_{i^*t},x_{i^*t-1},z_{i^*t},\widehat{\Delta\epsilon}_{i^*t})$,
with the \emph{unrestricted} residual $\widehat{\Delta\epsilon}_{it}=\Delta y_{it}-\Delta x_{it}'\hat\beta-\ind_{it}(\hat\gamma)'X_{it}\hat\delta$;
\item generate $\Delta y^*_{it}=\Delta x^{*\prime}_{it}\beta^*_0+\ind^*_{it}(\gamma^*_0)'X^*_{it}\delta^*_0+\widehat{\Delta\epsilon}^*_{it}$;
\item $g^*_{it}(\theta)=z^*_{it}(\Delta y^*_{it}-\Delta x^{*\prime}_{it}\beta-\ind^*_{it}(\gamma)'X^*_{it}\delta)$;
\item recentre: $\bar g^*_n(\theta)=\frac1n\sum_i\bigl(g^*_i(\theta)-\bar g_n(\hat\theta)\bigr)$, with
$\bar g_n(\hat\theta)$ the sample moment of the unrestricted residuals;
\item $\hat\theta^*_{(1)}=\arg\min\bar g^{*\prime}_n\bar g^*_n$ (identity weight) and
$W^*_n=\bigl(\frac1n\sum g^*_ig^{*\prime}_i-\bar g^*\bar g^{*\prime}\bigr)^{-1}$ at $\hat\theta^*_{(1)}$;
\item $\hat Q^*_n(\theta)=\bar g^*_n(\theta)'W^*_n\bar g^*_n(\theta)$ and the bootstrap statistic.
\end{enumerate}
Three choices of $\theta^*_0$ give three procedures. The grid bootstrap of the threshold takes
$\theta^*_0=(\hat\alpha(\gamma_\ell)',\gamma_\ell)'$ (the restricted estimate at the candidate) and
$D^*_n(\gamma_\ell)=n\bigl(\min_\alpha\hat Q^*_n(\alpha,\gamma_\ell)-\min_\theta\hat Q^*_n(\theta)\bigr)$;
the set is $\{\gamma\in\Gamma_n:D_n(\gamma)\le\hat F^{*-1}(1-\tau;D^*_n(\gamma))\}$ (eq.~7), which is never empty because
$D_n(\hat\gamma)=0$. Theorem~5: at $\gamma=\gamma_0$, $D^*_n\to Z_0^2$ with $Z_0=\max(0,Z^*_0)$ if the model is continuous and
$\to\chi^2_1$ if it is not, so the set is valid whether or not the function is continuous. The coefficient
bootstrap takes $\theta^*_0=w_n\hat\theta+(1-w_n)\tilde\theta$ with $w_n=\min\{T_n/(\hat Cn^{1/4}),1\}$ and the
continuity-restricted estimator $\tilde\theta$ (eq.~8), so that the bootstrap Jacobian degenerates fast enough
under continuity; the intervals are eq.~(9) (percentile) or eq.~(10) (symmetric), Theorem~6, pointwise only. The continuity test
takes $\theta^*_0=\tilde\theta$ and $T^*_n=n(\min_{\Theta_c}\hat Q^*_n-\min_\Theta\hat Q^*_n)$ (Theorem~7).

\paragraph{Comparison.}
\begin{longtable}{@{}p{3.1cm}p{6.1cm}p{6.8cm}@{}}
\toprule
Step & Gong and Seo, Algorithm 1 & \code{boot_type="wild"} (default) \\
\midrule
\endhead
Resampling & units drawn with replacement; $x,z$ and the residual move together & one Mammen weight per unit multiplies the residuals; $x$ and $z$ keep their sample values \\
Data-generating process & $\theta^*_0$ and \emph{unrestricted} residuals & restricted fit at the candidate and \emph{restricted} residuals; $\E^*[\eta]=0$, so the null holds without recentring \\
Recentring & subtract $\bar g_n(\hat\theta)$ & not needed (restricted residuals, mean-zero weights) \\
Weight in the bootstrap criterion & re-estimated in every draw from a first step with the identity & the sample weight, held fixed ($W_2$, or $W_1$ after the fallback) \\
Validity & Theorem 5 (continuous or not) & not covered by the theorem; coverage is examined by simulation (Nguyen and Lai 2026) \\
\bottomrule
\end{longtable}

\noindent \code{boot_type="unit"} follows the structure instead: whole units resampled (as count vectors),
restricted coefficients $\hat\theta(\gamma_\ell)$ with the \emph{unrestricted} residuals at $(\hat\theta,\hat\gamma)$, moments recentred at the sample moment
of those residuals, and $W^*_2$ re-estimated in every draw. One difference is deliberate in the Stata command and
copied: the first step of each draw uses the sample first-step weight $W_1$ where the paper uses the identity. It
is available for \code{method="fd"} without \code{kink} or \code{td}, and is much slower. An unbalanced effective
stack takes the command outside the common-$T$ setting of the theory (Section~\ref{note:map} lists the warning
that Stata prints and the package does not).

\paragraph{The coefficient bootstrap is not the paper's.} The legacy \code{coef_boot} generates data from
$\bX(\hat\gamma)\hat\theta$ with reweighted residuals, i.e.\ $\theta^*_0=\hat\theta$, which is the choice $w_n=1$ of
eq.~(8). That is the right choice when the model is discontinuous; under continuity the paper shows that a
bootstrap built on $\hat\theta$ is not valid, which is why it proposes $w_n$. The command therefore calls it a
legacy compatibility feature, \code{coef_boot="gs"} is refused, and the package says so in a note when it is used.
The symmetric interval is eq.~(10) with $\theta^*_0=\hat\theta$.

\paragraph{What is open.} The comparison above is a reading, not a proof of agreement. The power of the
rejection of shifted thresholds in the package's Monte Carlo experiments on the designs of Gong and Seo is lower than
their Table~2 while the coverage of the threshold agrees, and with the self-exciting designs of Seo and Shin
(2016) the dispersion of the estimated threshold is larger than in their tables, although the extreme estimates are
identical to Stata's, so they come from the estimator. The cause is not established.

\subsection{\texttt{xthenreg}}\label{note:xthenreg}

\paragraph{Estimator.} The balanced-panel estimator of Seo, Kim and Kim (2019) is the first-difference special
case of Section~\ref{note:est}: equation periods $t=2,\dots,T$, a block of instruments per period
(constant, $\diff x_t$ of the plain regressors, the levels $y_1,\dots,y_{t-2}$, the levels of the endogenous
regressors up to $t-2$, the external instruments), the coefficient order $[\beta\,|\,\delta_1\,|\,\delta]$ for the
jump and $[\beta\,|\,\kappa]$ for the kink, the first-step MA(1) weight built block by block
($W_1=(2S_0-S_1-S_1')^{-1}$ with $S_0$ block diagonal, $(1/N)Z_p'Z_p$, and $S_1$ block upper diagonal, $(1/N)Z_{p-1}'Z_p$), and the
second-step weight from the centred covariance of the unit moments. The grid is uniform between two
\code{mm_quantile} bounds at \code{trim_rate}/2 and $1-$\code{trim_rate}/2 over all the values of $q$ (default
\code{trim_rate} $0.4$) and $\hat\gamma$ is the first minimum of the second-step criterion; both steps
search the grid, the second step with the weight built at the first-step estimate.

\paragraph{Covariance.} Following Seo and Shin (2016, p.~173), with $G=(G_b,G_d,G_r)$ the derivatives of the
moments with respect to the linear coefficients, the regime coefficients and the threshold,
\begin{equation}\label{eq:xhcov}
\widehat V=\frac1N\,\bigl(G'\,\Omega^{-1}\,G\bigr)^{-1},
\end{equation}
with the derivative in $\gamma$ computed by the same Gaussian kernel as in (\ref{eq:xg}), bandwidth
$h=h_0\cdot1.06\cdot N^{-1/5}\cdot\mathrm{sd}(q)$ (\code{h_0(1.5)}), in the jump model, and in closed form for the kink.
The 95\,\% interval of every coefficient, the threshold included, is the estimate $\pm1.959963984540054\,\mathrm{se}$.

\paragraph{Linearity test.} The Wald statistic of the regime coefficients is computed at every grid point, each time
with its own two-step estimate and covariance, and the statistic is the supremum over the grid
(\code{e(sup_wald)} in the package). The p-value is from the fast bootstrap of the paper: the step-two moments of each
unit are multiplied by a standard normal weight ($\texttt{invnormal(uniform())}$), the one-step estimator is
re-evaluated with the fixed step-two weight at every grid point, and the Wald form reuses the covariance matrices of the
sample. The p-value is the \emph{share} of bootstrap suprema above the sample supremum, not the add-one rule of
\code{xtdpthresh}, so it can be exactly zero; a missing bootstrap statistic counts as exceeding.

\paragraph{Against \texttt{xtdpthresh}.} Under \code{method(fd)} the two commands use the same moment conditions and
weight matrices and, on the same grid, give the same estimates; the differences are in the trimming default, the
placement of $q$, the exogenous-regressor instruments (levels in \code{xthenreg}, transformed values in
\code{xtdpthresh}), the refinement of the kink, the absence of a standard error for $\hat\gamma$ and the GMM-distance
linearity test of \code{xtdpthresh}. Where the two designs coincide (cases c1, c4 and c5 of the validation) the package
checks that the two ports agree.

\subsection{Conventions of the Stata commands that are reproduced}\label{note:quirks}

Each row is a behaviour of the Stata code that the package copies so that the numbers agree. The ones marked
\code{STATA-QUIRK} in the source are the ones where the behaviour is surprising.

\begin{longtable}{@{}p{3.6cm}p{12.4cm}@{}}
\toprule
Convention & What is reproduced and why \\
\midrule
\endhead
\code{invsym} &
Reads only the upper triangle of its argument; a sweep with the pivot on the largest remaining diagonal element; a pivot
is dropped when it is at most $10^{-9}$ times the \emph{original} diagonal element, at most $9\times10^{-20}$, or negative,
and its row and column of the result are zero (a generalised inverse). Measured on Stata 17 MP. For a well-conditioned
positive definite matrix no pivot can be dropped in any order, so a Cholesky factorisation replaces the sweep. \\
\code{cholinv} &
Reads only the lower triangle; missing if the factorisation fails or the reciprocal condition number is below $10^{-20}$. \\
\code{syminv} gate &
\code{xtdpthresh} admits an inverse only after Jacobi equilibration, a Cholesky factorisation, a condition number of at most
$10^{12}$ and a check that $CC^{-1}\approx I$. A matrix that fails is a candidate that is not admissible. Validated on
23 matrices against the Mata source run unchanged. \\
Tolerances of ties &
Equal criteria (relative $10^{-12}$ for profiles, $10^{-10}$ for the refinement check) go to the smaller threshold
(\code{objtol}). \\
$W_2$ built once &
The second-step weight is built at the first-step threshold and held fixed over the grid of the second step. \\
Equal splits, exact zero &
Two candidates with the same split of the observed values of $q$ give the same regressors; the statistic of the estimated
threshold must be exactly zero, so the moments of the initial-grid entry with the same split are reused instead of
summing the arrays in another order (a gap of the order of rounding would otherwise remain). \\
Order statistic &
The critical value is the order statistic of rank $\lceil p(B+1)\rceil$, which is the add-one p-value rule; infinite when the
rank exceeds $B$ (no draw can reject). \\
Random stream &
Mersenne Twister 64, $u=((x\gg11)|1)2^{-53}$, row-major fill, component seeds $(s+\text{offset})\bmod2^{31}$, one draw per
contributing unit (a unit without equation takes no draw, so adding such a unit leaves the draws unchanged). \\
Unit bootstrap draws &
A draw whose profile has no usable minimum is skipped, not replaced (\code{STATA-QUIRK}); the missing ones stay missing in
the result. \\
\code{dlt_len} in \code{xthenreg} &
With a kink the Wald block that is tested has \emph{two} entries (the last linear coefficient and the kink
coefficient) instead of the kink coefficient alone (\code{STATA-QUIRK}); kept so that the statistic equals Stata's. \\
Signs of $G$ &
With a kink $G_d$ has the positive sign (unlike the jump model) and $G_r$ uses the indicator of the regime above the
threshold, as in \code{cov_mat_con_kink} (\code{STATA-QUIRK}). \\
\code{exogenous()} of \code{xthenreg} &
An \code{exogenous()} regressor enters with its levels up to $t$, two more lags than an \code{endogenous()} one
(\code{STATA-QUIRK}); hence 87 moment conditions where \code{xtdpthresh} has 60 instruments for the same data. \\
\code{boots_p} &
The share of bootstrap suprema above the sample supremum, a missing one counting as exceeding, as a missing value does in
Mata; $-1$ when the test is not run. \\
$h$ and $s_q$ &
$s_q$ is the sample standard deviation over \emph{all} the level observations and $n$ the units with an equation; the
1.959963984540054 of the interval is the constant typed in the \code{xthenreg} source. \\
AR on first differences &
The serial-correlation tests use first-difference residuals under both transformations. \\
Time-effect zeroing &
A column that the projection reduces to rounding noise (below $10^{-10}$ of its size) is set to exactly zero, as Stata does. \\
\code{if}, missing values &
A missing value in the condition counts as \emph{true} in Stata; in the package it counts as false, and the rows with a missing
variable are removed afterwards by the missing-value check, so the sample is the same. \\
Seed (the package adds) &
When the bootstrap is on and no seed is given, the wrapper draws one in $[1,2^{31}-1)$ and records it (\code{res.seed}, the
settings log, the JSON). Stata draws from the current generator and does not report it. \\
Precision &
All values are \code{float64}; Stata's operators on \code{float} variables are recomputed in double precision in the command
(\code{_xdpt_tsdouble}), which a double-precision array makes unnecessary. \\
\bottomrule
\end{longtable}

\subsection{Validation journal}\label{note:valid}

\paragraph{Rule.} Every number the package reports must equal Stata 17 MP to four decimals: absolute difference at most
$5\times10^{-5}$ \emph{or} relative difference at most $10^{-6}$. The comparison is automated and runs without Stata.

\paragraph{Method.} The real commands (\code{xthenreg}, \code{xtdpthresh} 0.9.38) are run in batch from PowerShell; a Mata
routine exports each object at 15 significant digits; the results are frozen in \code{tests/golden/} (one file per
component) with the versions of Stata and of the commands. The test suite compares the package with them case by case.
The random numbers of Stata are reproduced, so a seeded bootstrap draws the very weights of the command and its
p-values and sets are compared directly, not through their distribution. Each case is described once and drives both the
Stata run and the Python run.

\begin{longtable}{@{}p{4.3cm}p{2.2cm}p{9.5cm}@{}}
\toprule
Object & Cases & Result \\
\midrule
\endhead
Random numbers & 5 seeds & $10\,000$ draws each, identical; $40$ bootstrap suprema of each of 9 \code{xthenreg} cases equal to a relative $2\times10^{-10}$ \\
\code{invsym}, \code{cholinv}, \code{syminv} & measured; 23 matrices & thresholds measured on Stata; \code{syminv} equal to the Mata source run unchanged \\
\code{xthenreg}: $L$, $y$, $q$, $x$, $Z$ & 12 & equal to $10^{-12}$ running the Mata source as is (step~1 of \code{maincode_0719}) \\
\code{xthenreg}: $\hat\theta$, $\hat\gamma$, $V$, CI, $e()$ & c1--c8 and variants & $b$ within $3\times10^{-11}$ worst case on c1--c6; \code{h_0(3)} and \code{h_0(0.7)} (c7, c8) through the public function; grid $\times$ trimming variants (12 jump, 3 kink, 4 on $T=15$ with $L=325$) \\
\code{xthenreg}: sup-Wald, $p$ & 9; 5 & statistic to a relative $5\times10^{-10}$; the p-values of \code{boost(100)} equal on 5 cases (static and dynamic, jump and kink) \\
\code{xtdpthresh}: sample selection & 36 & \code{e(N)}, \code{e(N_stack)}, \code{e(N_units)} and the exact set of rows of \code{e(sample)} equal; return codes of the refusals equal \\
Instruments $Z$, $X$ & 28 + 30 + 15 & default, \code{maxlag}, \code{collapse}, \code{iv()}, duplicates, dependence, common variable; 15 cases of forward deviations \\
One step, two steps & 28 & \code{e(b)}, \code{e(obj)}, \code{e(V_cond)}; the one-step fallback is reproduced \\
Grid and search & 29; 5; 3 & \code{e(grid_admitted)}, \code{e(q_lo)}, \code{e(q_hi)}; kink refinement (\code{e(kink_refined)}); \code{refine} (iterations and points added) \\
Whole commands, full panel & 14 & the examples of the help file and variants on 565 firms: \code{e(b)}, \code{e(gamma)}, \code{e(obj)}, \code{e(N)}, \code{e(N_iv)} \\
Unbalanced panels, FOD, $td$, history & 5 + 5; 7 + 8; 4 & interior gaps, missing values, short units, alternate years; \code{td} under FD and FOD; \code{history(panel|sample)} \\
Variances & 45; 23 & \code{e(V)}, \code{e(V_cond)}, \code{e(V_upper)}, \code{e(b_upper)}, jump and kink, FD and FOD, \code{td}, unbalanced; \code{vce(windmeijer)}; \code{bwscale} (0.8, 3), \code{nocenter} \\
Tests & 45 & \code{e(hansen)}, \code{e(hansen_df)}, \code{e(hansen_p)}, \code{e(ar1)}, \code{e(ar2)} and their p-values \\
\code{predict} & 11 & residuals, \code{xb}, AR residuals, regime, \code{if} \\
Bootstrap & 62 Stata cases, 287 checks & \code{e(pval_lin)} (56), \code{e(ci_grid)}, \code{e(gamma_lo)}, \code{e(gamma_hi)}, \code{e(ci_nseg)}, \code{e(ci_segments)}; continuity (8); \code{citest} (8); coefficient bootstrap (10); \code{boottype(unit)}; refusals \\
Reporting & 146 checks & the tables repeat the validated $b$, $V$, \code{ci_grid}, \code{e(b_upper)} exactly \\
Hostile inputs & 19 & minimal panels, a constant, binary or rare threshold, an empty regime, duplicated regressors, fewer units than instruments, a null outcome: the same return code as Stata on 19 of 19\\
\bottomrule
\end{longtable}

\paragraph{Equivalence of the fast kernels.} Each Numba kernel has a NumPy reference that is kept; the tests compare them to
$10^{-12}$ (several seeds, ties between $q$ and the grid, a missing $q$, degenerate grids, one against four threads bit for
bit, and with the compiler switched off). The optimisations of the bootstrap were accepted only when the reference
cases remained identical.

\paragraph{Suite.} About 2\,170 tests (unit, golden, integration), line and branch coverage about 91.5\,\% (the compiled bodies
cannot be traced by the coverage tool); the code of the README and of this guide is executed by the tests.

\paragraph{What is not compared with Stata.} The statistical behaviour of the package's own seeds (size and power of the
linearity test, coverage of the confidence set over many panels) is measured by Monte Carlo and compared with the
papers, not with Stata, and has the open points of Section~\ref{note:gs}. Timings are measured on the same machine one run
after the other, because the speed of a machine varies by up to a factor of two from one moment to the next. The statements
of the guide about interpretation come from the Stata help files, the code and the execution of the package, and are not
a second validation.

\subsection{Where every Stata function went}\label{note:map}

The inventory is generated by \code{tools/ported_functions.py} from the \texttt{\# Stata: function (file.ado:line)} comments of
the source and the function ranges of \code{validation/stata_functions.json} (a manifest with the SHA-256 of the three
\code{.ado} files); a test checks that every citation points inside a function. A function counts as \emph{ported} when
at least one comment points inside it. At the time of writing, 86 of the 101 functions of the three files are cited; the generated
tables below list the cited ones and where they live. The 15 that are not cited fall in three groups.

\begin{longtable}{@{}p{3.7cm}p{3.0cm}p{9.3cm}@{}}
\toprule
Stata function & Status & Reason \\
\midrule
\endhead
\code{xdpt2_append_nested_level}, \code{xdpt2_nested_midpoints}, \code{xdpt2_profile_stability} &
retired &
Reachable only when \code{search_mode == "adaptive"}, which the command refuses since 0.9.26 (\code{search_mode}, \code{search_max} and
\code{search_tol} are retired). The package evaluates the grid it is asked for; the nested grids $g\to2g-1\to4g-3$ are reached by choosing
\code{grid=199} or \code{397}. \\
\code{xdpt2_build_gamma_cache} &
dead code &
No caller in the command; the variant \code{xdpt2_build_gamma_cache_t} is the one in use and is ported (the cache of the candidates). \\
\code{xdpt2_tpl_build}, \code{xdpt2_tpl_dW}, \code{xdpt2_tpl_release}, \code{xdpt2_unit_cfg} &
replaced &
The Mata record of the stacked-equation template, its release and the settings read once per stack. The package keeps the
transformation as the sparse matrix $\mathcal T$ (Section~\ref{note:model}) and the instrument settings as a dataclass, so none of this
bookkeeping exists; the quantities they produce are the ones compared with Stata. \\
\code{xdpt2_hasmiss} &
not needed &
Tests for a missing value in a block; the package uses \code{np.isfinite}. \\
\code{_xdpt_tsdouble} &
not needed &
Recomputes the time-series operators in double precision (\code{tsrevar} keeps the type of the source); all values are \code{float64}. \\
\code{xdpt2_rowsig2} (in \code{xtdpthresh.ado} and in \code{xtdpthresh_p.ado}), \code{xdpt2_p_fill} of \code{xtdpthresh_p.ado} &
replaced &
A checksum of the data and a keyed store that let \code{predict} detect a changed dataset or a cleared memory (return codes 459 and
498). The result object holds its data and \code{predict(kind, frame)} checks the keys of the frame it is given; the copy of
\code{xdpt2_p_fill} in \code{xtdpthresh.ado} is ported. \\
\code{xdpt2_is_strongly_balanced} &
message not ported &
Used to print, when \code{boottype(unit)} runs on an unbalanced effective stack, ``this is an experimental extension beyond the
common-$T$ theory'', and to store a balance flag. The package does not print that warning; the estimates are unaffected. \\
\code{xdpt2_near_common} &
message not ported &
Names, in the error message of a rank failure, the variables that are common to all units only up to rounding. The package raises
the same error with the same return code but without that hint. \\
\bottomrule
\end{longtable}

\noindent The 19 functions that were not cited before this note was written are now cited where the package implements their role
(for example \code{xdpt2_gsum_by_unit} by the unit-moment kernel, \code{xdpt2_kink_joint_V} and
\code{xdpt2_kink_joint_wind} by the joint variance, \code{xdpt2_fast_gmm_boot} by the batched bootstrap criterion,
\code{FD_con}, \code{cov_con_b} and \code{cov_con_kink_b} of \code{xthenreg} by the first differences and the Wald
covariance of the linearity test).

Generated by \texttt{tools/ported\_functions.py\ inventory} from the
\texttt{\#\ Stata:} comments of \texttt{src/} and the function ranges of
\texttt{validation/stata\_functions.json}. A function is \emph{ported}
when at least one comment of the package points inside it.

\subsubsection{\texorpdfstring{\texttt{xtdpthresh.ado} --- 63 of 76
functions
referenced}{xtdpthresh.ado --- 63 of 76 functions referenced}}\label{xtdpthresh.ado-63-of-76-functions-referenced}

{\def\LTcaptype{none} % do not increment counter
\begin{longtable}[]{@{}
  >{\raggedright\arraybackslash}p{(\linewidth - 4\tabcolsep) * \real{0.2700}}
  >{\raggedright\arraybackslash}p{(\linewidth - 4\tabcolsep) * \real{0.0900}}
  >{\raggedright\arraybackslash}p{(\linewidth - 4\tabcolsep) * \real{0.6400}}@{}}
\toprule\noalign{}
\begin{minipage}[b]{\linewidth}\raggedright
Stata function
\end{minipage} & \begin{minipage}[b]{\linewidth}\raggedright
Lines
\end{minipage} & \begin{minipage}[b]{\linewidth}\raggedright
Ported in
\end{minipage} \\
\midrule\noalign{}
\endhead
\bottomrule\noalign{}
\endlastfoot
\texttt{\_xdpt\_coeftab} & 2319-2367 &
\texttt{extras/reporting/summary.py} \\
\texttt{\_xdpt\_tsterm} & 2391-2466 & \texttt{data/operators.py} \\
\texttt{\_xdpt\_upper} & 2257-2318 &
\texttt{models/xtdpthresh/coefficients.py} \\
\texttt{xdpt2\_append\_candidates} & 5275-5358 &
\texttt{models/xtdpthresh/search.py} \\
\texttt{xdpt2\_ar\_full} & 8092-8260 &
\texttt{models/xtdpthresh/diagnostics.py} \\
\texttt{xdpt2\_ar\_xg} & 5149-5190 &
\texttt{models/xtdpthresh/stack.py} \\
\texttt{xdpt2\_build\_W\_ma1} & 3808-3879 & \texttt{core/weights.py},
\texttt{models/xtdpthresh/stack.py} \\
\texttt{xdpt2\_build\_cluster\_omega} & 4001-4011 &
\texttt{models/xtdpthresh/estimator.py} \\
\texttt{xdpt2\_build\_cluster\_omega\_c} & 4012-4047 &
\texttt{core/gmm.py} \\
\texttt{xdpt2\_build\_gamma\_cache\_t} & 4383-4619 &
\texttt{core/grids.py}, \texttt{inference/bootstrap/objective.py},
\texttt{models/xtdpthresh/search.py},
\texttt{models/xtdpthresh/stack.py} \\
\texttt{xdpt2\_build\_units} & 2587-2689 &
\texttt{models/xtdpthresh/sample.py} \\
\texttt{xdpt2\_cache\_w2} & 6647-6666 &
\texttt{inference/bootstrap/objective.py} \\
\texttt{xdpt2\_coef\_bootstrap} & 6357-6633 &
\texttt{models/xtdpthresh/bootstrap.py} \\
\texttt{xdpt2\_component\_seed} & 4796-4806 &
\texttt{inference/bootstrap/draws.py},
\texttt{inference/bootstrap/rng.py} \\
\texttt{xdpt2\_const\_within} & 8261-8304 &
\texttt{models/xtdpthresh/sample.py} \\
\texttt{xdpt2\_continuity\_test} & 7431-7767 &
\texttt{models/xtdpthresh/bootstrap.py} \\
\texttt{xdpt2\_crit\_orderstat} & 4839-4888 &
\texttt{inference/bootstrap/objective.py} \\
\texttt{xdpt2\_demean\_bytime} & 3142-3176 &
\texttt{transforms/td.py} \\
\texttt{xdpt2\_dense\_uid} & 4740-4759 &
\texttt{models/xtdpthresh/estimator.py} \\
\texttt{xdpt2\_drop\_cellconst} & 3304-3399 &
\texttt{instruments/cleaning.py} \\
\texttt{xdpt2\_fast\_gmm\_boot} & 4620-4667 &
\texttt{inference/bootstrap/objective.py} \\
\texttt{xdpt2\_fast\_gmm\_boot\_w} & 6667-6696 &
\texttt{inference/bootstrap/objective.py} \\
\texttt{xdpt2\_fast\_obj\_split\_list} & 4668-4739 &
\texttt{inference/bootstrap/objective.py} \\
\texttt{xdpt2\_find\_t} & 2703-2720 & \texttt{core/calendar.py} \\
\texttt{xdpt2\_gmm\_sandwich} & 4048-4099 & \texttt{core/gmm.py} \\
\texttt{xdpt2\_grid\_bootstrap} & 6697-7430 &
\texttt{models/xtdpthresh/bootstrap.py},
\texttt{models/xtdpthresh/unit\_bootstrap.py} \\
\texttt{xdpt2\_grid\_search} & 5630-6356 &
\texttt{models/xtdpthresh/estimator.py},
\texttt{models/xtdpthresh/search.py} \\
\texttt{xdpt2\_gsum\_by\_unit} & 3957-4000 &
\texttt{models/xtdpthresh/estimator.py} \\
\texttt{xdpt2\_hansen\_j} & 8048-8091 &
\texttt{models/xtdpthresh/diagnostics.py} \\
\texttt{xdpt2\_indep\_cols} & 3177-3243 &
\texttt{instruments/cleaning.py} \\
\texttt{xdpt2\_jump\_xg} & 5067-5120 &
\texttt{models/xtdpthresh/stack.py} \\
\texttt{xdpt2\_jump\_xg\_tp} & 5035-5066 &
\texttt{models/xtdpthresh/stack.py} \\
\texttt{xdpt2\_kink\_apply\_joint} & 5237-5256 &
\texttt{models/xtdpthresh/variance.py} \\
\texttt{xdpt2\_kink\_joint\_V} & 5191-5209 &
\texttt{models/xtdpthresh/variance.py} \\
\texttt{xdpt2\_kink\_joint\_wind} & 5210-5236 &
\texttt{models/xtdpthresh/variance.py} \\
\texttt{xdpt2\_kink\_refine} & 5516-5629 &
\texttt{models/xtdpthresh/search.py} \\
\texttt{xdpt2\_kink\_xg} & 4954-4986 &
\texttt{models/xtdpthresh/stack.py} \\
\texttt{xdpt2\_linearity\_test} & 7768-8047 &
\texttt{models/xtdpthresh/bootstrap.py} \\
\texttt{xdpt2\_mammen\_draw} & 4783-4795 &
\texttt{inference/bootstrap/draws.py} \\
\texttt{xdpt2\_n\_switch} & 4274-4308 &
\texttt{models/xtdpthresh/stack.py} \\
\texttt{xdpt2\_nodup\_cols} & 3244-3303 &
\texttt{instruments/cleaning.py} \\
\texttt{xdpt2\_objtol} & 2497-2515 & \texttt{core/linalg.py} \\
\texttt{xdpt2\_p\_fill} & 8349-8434 &
\texttt{models/xtdpthresh/predict.py} \\
\texttt{xdpt2\_profile\_is\_flat} & 5456-5469 &
\texttt{models/xtdpthresh/search.py} \\
\texttt{xdpt2\_profile\_update} & 5386-5455 &
\texttt{models/xtdpthresh/search.py} \\
\texttt{xdpt2\_psd\_ok} & 2554-2586 & \texttt{core/linalg.py} \\
\texttt{xdpt2\_q\_at\_rows} & 2721-2743 &
\texttt{models/xtdpthresh/search.py} \\
\texttt{xdpt2\_q\_nvals} & 5121-5148 &
\texttt{models/xtdpthresh/variance.py} \\
\texttt{xdpt2\_q\_support} & 2744-2790 &
\texttt{models/xtdpthresh/search.py} \\
\texttt{xdpt2\_quantile} & 4807-4838 & \texttt{core/quantiles.py} \\
\texttt{xdpt2\_quantile\_grid} & 2467-2479 & \texttt{core/grids.py} \\
\texttt{xdpt2\_rangen} & 2480-2496 & \texttt{core/quantiles.py} \\
\texttt{xdpt2\_solve\_gmm\_1step} & 3880-3916 & \texttt{core/gmm.py} \\
\texttt{xdpt2\_solve\_gmm\_1step\_pre} & 3917-3956 &
\texttt{core/gmm.py} \\
\texttt{xdpt2\_stack\_at\_gamma} & 3452-3807 &
\texttt{instruments/cleaning.py}, \texttt{models/xtdpthresh/stack.py},
\texttt{transforms/td.py} \\
\texttt{xdpt2\_syminv} & 2516-2553 & \texttt{core/linalg.py} \\
\texttt{xdpt2\_tpl\_trans} & 4987-5034 &
\texttt{models/xtdpthresh/stack.py} \\
\texttt{xdpt2\_tpos} & 2690-2702 & \texttt{core/calendar.py} \\
\texttt{xdpt2\_transform\_unit} & 2801-3141 &
\texttt{instruments/builder.py}, \texttt{models/xtdpthresh/stack.py},
\texttt{transforms/fd.py}, \texttt{transforms/fod.py} \\
\texttt{xdpt2\_w2\_solvable} & 6634-6646 &
\texttt{inference/bootstrap/objective.py} \\
\texttt{xdpt2\_windmeijer} & 4889-4953 &
\texttt{models/xtdpthresh/variance.py} \\
\texttt{xtdpthresh} & 5-2256 & \texttt{models/xtdpthresh/bootstrap.py},
\texttt{models/xtdpthresh/coefficients.py},
\texttt{models/xtdpthresh/design.py}, \texttt{models/xtdpthresh/fit.py},
\texttt{models/xtdpthresh/options.py},
\texttt{models/xtdpthresh/sample.py} \\
\texttt{xtdpthresh\_run} & 8435-11202 & \texttt{core/grids.py},
\texttt{instruments/builder.py},
\texttt{models/xtdpthresh/bootstrap.py},
\texttt{models/xtdpthresh/diagnostics.py},
\texttt{models/xtdpthresh/fit.py}, \texttt{models/xtdpthresh/search.py},
\texttt{models/xtdpthresh/variance.py} \\
\end{longtable}
}

Not referenced by any comment (13): \texttt{\_xdpt\_tsdouble},
\texttt{xdpt2\_append\_nested\_level},
\texttt{xdpt2\_build\_gamma\_cache}, \texttt{xdpt2\_hasmiss},
\texttt{xdpt2\_is\_strongly\_balanced}, \texttt{xdpt2\_near\_common},
\texttt{xdpt2\_nested\_midpoints}, \texttt{xdpt2\_profile\_stability},
\texttt{xdpt2\_rowsig2}, \texttt{xdpt2\_tpl\_build},
\texttt{xdpt2\_tpl\_dW}, \texttt{xdpt2\_tpl\_release},
\texttt{xdpt2\_unit\_cfg}

\subsubsection{\texorpdfstring{\texttt{xtdpthresh\_p.ado} --- 1 of 3
functions
referenced}{xtdpthresh\_p.ado --- 1 of 3 functions referenced}}\label{xtdpthresh_p.ado-1-of-3-functions-referenced}

{\def\LTcaptype{none} % do not increment counter
\begin{longtable}[]{@{}
  >{\raggedright\arraybackslash}p{(\linewidth - 4\tabcolsep) * \real{0.2700}}
  >{\raggedright\arraybackslash}p{(\linewidth - 4\tabcolsep) * \real{0.1500}}
  >{\raggedright\arraybackslash}p{(\linewidth - 4\tabcolsep) * \real{0.5800}}@{}}
\toprule\noalign{}
\begin{minipage}[b]{\linewidth}\raggedright
Stata function
\end{minipage} & \begin{minipage}[b]{\linewidth}\raggedright
Lines
\end{minipage} & \begin{minipage}[b]{\linewidth}\raggedright
Ported in
\end{minipage} \\
\midrule\noalign{}
\endhead
\bottomrule\noalign{}
\endlastfoot
\texttt{xtdpthresh\_p} & 49-218 &
\texttt{models/xtdpthresh/predict.py} \\
\end{longtable}
}

Not referenced by any comment (2): \texttt{xdpt2\_p\_fill},
\texttt{xdpt2\_rowsig2}

\subsubsection{\texorpdfstring{\texttt{xthenreg.ado} --- 22 of 22
functions
referenced}{xthenreg.ado --- 22 of 22 functions referenced}}\label{xthenreg.ado-22-of-22-functions-referenced}

{\def\LTcaptype{none} % do not increment counter
\begin{longtable}[]{@{}
  >{\raggedright\arraybackslash}p{(\linewidth - 4\tabcolsep) * \real{0.1900}}
  >{\raggedright\arraybackslash}p{(\linewidth - 4\tabcolsep) * \real{0.0900}}
  >{\raggedright\arraybackslash}p{(\linewidth - 4\tabcolsep) * \real{0.7200}}@{}}
\toprule\noalign{}
\begin{minipage}[b]{\linewidth}\raggedright
Stata function
\end{minipage} & \begin{minipage}[b]{\linewidth}\raggedright
Lines
\end{minipage} & \begin{minipage}[b]{\linewidth}\raggedright
Ported in
\end{minipage} \\
\midrule\noalign{}
\endhead
\bottomrule\noalign{}
\endlastfoot
\texttt{FD\_con} & 885-895 & \texttt{models/xthenreg/moments.py} \\
\texttt{GMM\_W\_n\_2\_con} & 1293-1316 & \texttt{core/gmm.py},
\texttt{core/weights.py} \\
\texttt{GMM\_W\_n\_con} & 1110-1174 & \texttt{core/weights.py} \\
\texttt{GMM\_W\_n\_static} & 1349-1414 & \texttt{core/weights.py} \\
\texttt{GMM\_cal} & 1175-1207 & \texttt{core/gmm.py},
\texttt{models/xthenreg/moments.py} \\
\texttt{GMM\_cal\_kink} & 1415-1453 &
\texttt{models/xthenreg/moments.py} \\
\texttt{GMM\_var\_con} & 1014-1109 &
\texttt{models/xthenreg/moments.py} \\
\texttt{Z\_mat\_IV\_add} & 1454-1531 &
\texttt{models/xthenreg/design.py} \\
\texttt{Z\_mat\_con} & 896-1013 & \texttt{models/xthenreg/design.py} \\
\texttt{cov\_con\_b} & 2034-2079 &
\texttt{models/xthenreg/linearity.py} \\
\texttt{cov\_con\_kink\_b} & 2080-2117 &
\texttt{models/xthenreg/linearity.py} \\
\texttt{cov\_mat\_con} & 1532-1593 &
\texttt{models/xthenreg/covariance.py} \\
\texttt{cov\_mat\_con\_kink} & 1594-1635 &
\texttt{models/xthenreg/covariance.py} \\
\texttt{ep\_cal} & 1208-1249 & \texttt{models/xthenreg/moments.py} \\
\texttt{ep\_cal\_kink} & 1250-1292 &
\texttt{models/xthenreg/moments.py} \\
\texttt{grid\_con} & 870-884 & \texttt{core/grids.py} \\
\texttt{index\_non\_ms} & 1636-1682 &
\texttt{models/xthenreg/design.py} \\
\texttt{maincode\_0719} & 321-869 & \texttt{models/xthenreg/design.py},
\texttt{models/xthenreg/estimate.py},
\texttt{models/xthenreg/estimator.py},
\texttt{models/xthenreg/linearity.py} \\
\texttt{moment\_mat\_con} & 1317-1348 &
\texttt{models/xthenreg/design.py} \\
\texttt{sup\_wald\_cal} & 1683-1919 &
\texttt{models/xthenreg/linearity.py} \\
\texttt{sup\_wald\_fb} & 1920-2033 &
\texttt{models/xthenreg/linearity.py} \\
\texttt{xthenreg} & 21-320 & \texttt{models/xthenreg/estimate.py} \\
\end{longtable}
}

Not referenced by any comment (0): none


\subsection*{References}

\begin{description}\small
\item[Arellano, M., and S. Bond (1991).] Some tests of specification for panel data: Monte Carlo evidence and an application to
employment equations. \emph{Review of Economic Studies} 58: 277--297.
\item[Arellano, M., and O. Bover (1995).] Another look at the instrumental variable estimation of error-components models.
\emph{Journal of Econometrics} 68: 29--51.
\item[Gong, W., and M. H. Seo (2026).] Bootstraps for dynamic panel threshold models. \emph{Journal of Econometrics} 253: 106153.
\item[Hansen, B. E. (1999a).] Threshold effects in non-dynamic panels: estimation, testing, and inference.
\emph{Journal of Econometrics} 93: 345--368.
\item[Hansen, B. E. (1999b).] The grid bootstrap and the autoregressive model. \emph{Review of Economics and Statistics} 81: 594--607.
\item[Mammen, E. (1993).] Bootstrap and wild bootstrap for high dimensional linear models. \emph{Annals of Statistics} 21: 255--285.
\item[Nguyen, D. C., and N. D. Lai (2026).] xtdpthresh: dynamic panel threshold regression for unbalanced panels, with endogenous
regressors and continuity-robust inference. SSRN Working Paper 6619058.
\item[Roodman, D. (2009).] How to do xtabond2: an introduction to difference and system GMM in Stata. \emph{Stata Journal} 9: 86--136.
\item[Seo, M. H., S. Kim, and Y.-J. Kim (2019).] Estimation of dynamic panel threshold model using Stata. \emph{Stata Journal} 19: 685--697.
\item[Seo, M. H., and Y. Shin (2016).] Dynamic panels with threshold effect and endogeneity. \emph{Journal of Econometrics} 195: 169--186.
\item[Windmeijer, F. (2005).] A finite sample correction for the variance of linear efficient two-step GMM estimators.
\emph{Journal of Econometrics} 126: 25--51.
\end{description}

\end{document}
